\documentclass[12pt, a4paper]{article}
\usepackage[T1]{fontenc}
\usepackage{lmodern}
\usepackage{amsmath, amsthm, amssymb, mathtools}
\usepackage[margin=1in]{geometry}
\usepackage{xr-hyper}
\usepackage[colorlinks=true,linkcolor=blue,citecolor=blue]{hyperref}
\usepackage{cleveref}
\usepackage{graphicx, subcaption}
\usepackage{natbib}
\usepackage{booktabs}
\usepackage{enumerate}

% Code references: scaled monospace with line-breaking
\newcommand{\code}[1]{\texttt{\small #1}}

\newtheorem{theorem}{Theorem}[section]
\newtheorem{lemma}[theorem]{Lemma}
\newtheorem{proposition}[theorem]{Proposition}
\newtheorem{corollary}[theorem]{Corollary}
\newtheorem{definition}{Definition}[section]
\newtheorem{remark}{Remark}[section]
\newtheorem{conjecture}[theorem]{Conjecture}
\newtheorem{derivation}[theorem]{Derivation}
\newtheorem{observation}[theorem]{Observation}
\numberwithin{equation}{section}

\newcommand{\dd}{\mathrm{d}}
\newcommand{\seriestitle}{Why Rotating Fluids Make Polygons}


\externaldocument[I-]{../paper-1-mathematics/main}
\externaldocument[II-]{../paper-2-physics/main}
\externaldocument[III-]{../paper-3-gravity/main}
\externaldocument[IV-]{../paper-4-field-theory/main}
\externaldocument[V-]{../paper-5-s3-framework/main}
\externaldocument[SP-]{../supplement-proofs/main}

\title{\seriestitle\\[6pt]\large Paper VI: Cosmological Parameters}
\author{Gordon Speagle}
\date{}

\begin{document}
\maketitle

\begin{abstract}
We derive the cosmological constant, Hubble parameter, dark
energy/matter fractions, neutrino mass ratio, and baryon
asymmetry from the Havelock-Pell criterion
($N = 4$ and $N = 7$; Paper~IV), which determines
the cosmological polygon $N = 11 = 4 + 7$ via
Kaluza--Klein flux additivity on the derived $S^1$ fiber.
The $N = 11$ Born--Oppenheimer tunneling action fixes the
hyperbolic curvature radius~$\ell$, giving
$H_0 = 67.4$\,km/s/Mpc (Planck 2018 \citep{Planck2018}: $67.4 \pm 0.5$),
$\Omega_\Lambda = 68.3\%$,
$\Omega_{\mathrm{DM}} = 26.7\%$,
$\Omega_b = 4.97\%$ ($0.4\sigma$),
$\Delta m^2_{\mathrm{sol}}/\Delta m^2_{\mathrm{atm}} = 1/31.5$ ($3\%$),
and $\eta_B = 5.7 \times 10^{-10}$
($7\%$ match to the observed
$6.12 \times 10^{-10}$).
The theory predicts normal neutrino hierarchy,
$w = -1$ exactly, and no neutrinoless double beta decay.
\end{abstract}

\section{Introduction}
\label{sec:intro-cosmo}

Papers~I--IV developed the Havelock eigenvalue decomposition
$\lambda_m = C_1(\text{surface}) - m(N{-}m)/2$
as pure mathematics (Paper~I), applied it to vortex and
gravitational stability (Papers~II--III), and derived the
Standard Model gauge group and particle spectrum from the
$N = 7$ polygon sector on the Seifert manifold
$\mathbf{H}^2 \times_N S^1$ (Paper~IV).

This paper adds the \emph{cosmological sector}: the $N = 11$
polygon, selected by KK flux additivity
($N_{\mathrm{total}}/2 = N_{\mathrm{EW}}/2 + N_{\mathrm{grav}}/2
= 4/2 + 7/2 = 11/2$).
The $N = 11$ Born--Oppenheimer tunneling action sets the
curvature radius $\ell$ of the $\mathbf{H}^2$ base, thereby
determining the physical cosmological constant and the full
energy budget.

\paragraph{Notation and prerequisites.}
We use the Seifert manifold and field content established in
Paper~IV (\S\ref*{IV-sec:spacetime}--\S\ref*{IV-sec:fields}).
The central charge $c_N = 12\,b(N)$
(we write $c_N$ when a specific polygon order is intended;
$c_7 = 12\,b(7) = 51.57$, $c_{11} = 12\,b(11) = 126.56$),
the BO potential
$V_N(\rho) = \log(2\sinh\rho) + b(N) - f(m^*,N)$
(we write $V_N$ for the effective potential at polygon order~$N$),
and the Onsager--WDW correspondence (Paper~III) are
assumed throughout.
The instanton fugacity $K = e^{-2\pi k_{\mathrm{frac}}}$
and the WKB tunneling action $S_{\mathrm{BO}}(N)$ are
defined in Paper~IV, \S\ref*{IV-sec:couplings}.

\section{The cosmological constant}
\label{sec:lambda}

\begin{proposition}[Cosmological constant from the fiber]
\label{prop:lambda}
The effective three-dimensional cosmological constant is
\begin{equation}\label{eq:lambda}
  \Lambda_3 = \frac{N^2 - 16}{16}.
\end{equation}
$\Lambda_3 = 0$ at $N = 4$ (the gauge threshold, $j = 1$);
$\Lambda_3 < 0$ for $N = 3$ (AdS);
$\Lambda_3 > 0$ for $N \ge 5$ (de~Sitter).
At $N = 7$: $\Lambda_3 = (49{-}16)/16 = 2.0625$ (in $\ell^{-2}$ units). The physical $\Lambda_4 = \Lambda_3/\ell^2$, with $\ell$ from the $N = 11$ instanton (\S\ref{sec:cc-instanton}), gives $\Lambda_4 \sim 10^{-122}\,M_P^4$.
\end{proposition}

\begin{proof}
The cosmological constant is the competition between the
negative curvature of~$\mathbf{H}^2$ (which wants $\Lambda_3 < 0$)
and the positive flux energy on the $S^1$ fiber
(which wants $\Lambda_3 > 0$); at $N = 4$ these cancel exactly.

In the standard KK reduction of the metric~\eqref{IV-eq:metric-4d},
the fiber contribution to the effective potential is
$(1/4)|F|^2$, where $F = dA$ is the gauge field strength.
With quantized flux $\oint dA = \pi N$
(from~\eqref{IV-eq:flux}), the energy density on a unit-area
base is $(1/4)(\pi N / (2\pi))^2 = N^2/16$.
The base Gaussian curvature: $K_{\mathbf{H}^2} = -1$
(Ricci scalar $R = 2K = -2$; the magnitude $|K| = 1$ sets
the curvature scale).
The KK reduction of $\int R^{(4)}\sqrt{g^{(4)}}$
with metric $ds^2 = g_{ij}dx^idx^j + (d\varphi{+}A)^2$
gives $\int[R^{(3)} - \tfrac{1}{4}F_{ij}F^{ij}]\sqrt{g^{(3)}}$
(the $1/4$ is the standard Maxwell normalization from KK).
The 3D Einstein equations on the base
$G_{\mu\nu}^{(3)} + \Lambda_3\,g_{\mu\nu}^{(3)}
= 8\pi G_3\,T_{\mu\nu}^{(\mathrm{flux})}$
have $T_{\mu\nu}^{(\mathrm{flux})}
= (1/4)|F|^2\,g_{\mu\nu}^{(3)}$ (Maxwell stress from the
KK gauge field).
On a constant-curvature base with $K = -1$
(so $G^{(3)} = -K\,g = g$):
$\Lambda_3 = K + (1/4)|F|^2 = -1 + N^2/16$,
where $|F|^2 = (\pi N/(2\pi))^2 = N^2/4$ from the
quantized flux $\oint\!dA = \pi N$.
On $\mathbf{H}^2$ with $K = -1$:
$\Lambda_3 = -1 + N^2/16 = (N^2 - 16)/16$
(in units of $1/\ell^2$).
The $\mathbb{Z}_N$ cone points on the base do not
affect this formula: the KK reduction is a local
computation (the Einstein--Hilbert integrand is a
density), and the $\mathbb{Z}_N$ orbifold singularities
have zero measure.
The flux quantization $\oint dA = \pi N$ is
determined by the orbifold Euler class $e = N/2$,
which is a topological invariant of the Seifert
fibration (unchanged by resolving or smoothing the
cone points).
\end{proof}

\noindent
The physical four-dimensional cosmological constant is
$\Lambda_4 = \Lambda_3/\ell^2$,
where $\ell$ is the $\mathbf{H}^2$ curvature radius
(determined by the $N = 11$ tunneling action,
Theorem~\ref{thm:cc}).
The ratio $\Omega_\Lambda/\Omega_{\mathrm{DM}}$ is
dimensionless and $\ell$-independent.

\noindent
The two thresholds have a unified interpretation:
\begin{center}
\small
\begin{tabular}{lcll}
\toprule
$N$ & $\Lambda_3$ & $j_{\mathrm{crit}}$ & Physical meaning \\
\midrule
$4$ & $0$ & $j = 1$ & Gauge boson; $\Lambda_3 = 0$ boundary \\
$7$ & $33/16$ & $j = 2$ & Graviton; stability boundary $N_{\mathrm{crit}} = 7$ \\
\bottomrule
\end{tabular}
\end{center}


\section{The $N$ selection principle}
\label{sec:n-selection}

\begin{derivation}[$N$ selection from angular momentum quantization]
\label{thm:n-selection}
The three polygon numbers are determined by three structural
requirements (integer spin, stability, flux additivity):
\begin{enumerate}
\item[(a)] $N_{\mathrm{EW}} = 4$: the unique even polygon with
  integer critical spin ($j = 1$ from $f(2,4) = 2 = j(j{+}1)$).
  Equivalently: the $\Lambda_3 = 0$ threshold.
\item[(b)] $N_{\mathrm{grav}} = 7$: the unique odd polygon with
  integer critical spin in the stable range $N \le N_{\mathrm{crit}} = 7$
  ($j = 2$ from $f(3,7) = 6 = j(j{+}1)$).
\item[(c)] $N_{\mathrm{cosmo}} = 11 = N_{\mathrm{EW}} + N_{\mathrm{grav}}$:
  the two gauge sectors share a single $S^1$ fiber
  with KK fluxes $N_{\mathrm{EW}}/2 = 2$ and
  $N_{\mathrm{grav}}/2 = 7/2$.
  The total flux is the first Chern class of the
  Whitney sum of the two determinant line bundles:
  $c_1(L_{\mathrm{EW}} \oplus L_{\mathrm{grav}})
  = c_1(L_{\mathrm{EW}}) + c_1(L_{\mathrm{grav}})
  = 4 + 7 = 11$, giving $N = 11$.
  No geometric input enters, only the linearity of
  $\oint A$ on the fiber (exhaustive elimination below).
\end{enumerate}
\end{derivation}

\begin{proof}
\emph{Part (a).}
For even $N$: $f(N/2, N) = N^2/8$.
This is an integer only when $8 \mid N^2$, i.e.\ $4 \mid N$
(since $N^2/8 \in \mathbb{Z}$ requires $v_2(N^2) \ge 3$,
i.e.\ $v_2(N) \ge 2$, i.e.\ $4 \mid N$).
Writing $N = 4k$: $f(N/2, N) = 2k^2$.
The condition $j(j{+}1) = 2k^2$ with $j \in \mathbb{Z}$
requires $1 + 8k^2$ to be a perfect square.
At $k = 1$ ($N = 4$): $1 + 8 = 9 = 3^2$, giving $j = 1$.
The next solution is $k = 6$ ($N = 24$, $j = 8$).
So $N = 4$ is the unique even polygon with $j \le 7$.

\emph{Part (b).}
For odd~$N$: $f(\lfloor N/2 \rfloor, N) = (N^2{-}1)/8$.
Setting $y = 2j{+}1$, the integer-spin condition becomes
\begin{equation}\label{eq:pell}
  N^2 - 2y^2 = -1
  \qquad \text{(the negative Pell equation for $\sqrt{2}$)}.
\end{equation}
The solutions are generated by the fundamental unit
$\varepsilon = 1 + \sqrt{2}$ of $\mathbb{Z}[\sqrt{2}]$:
$(N_k, y_k) = $ the odd powers $\varepsilon^{2k-1}$.
Explicitly: $(1, 1)$, $\mathbf{(7, 5)}$, $(41, 29)$, $(239, 169)$, \ldots,
giving $j = 0, \mathbf{2}, 14, 84, \ldots$
The first non-trivial solution is $N = 7$, $j = 2$ (the graviton).
The next ($N = 41$) lies far above the flat-plane stability boundary
$N_{\mathrm{crit}} = 7$.

\emph{Part (c).}
The Seifert fiber $S^1$ is a circle bundle over $\mathbf{H}^2/\mathbb{Z}_N$,
classified by its degree (= first Chern class $c_1 = N$).
The two gauge sectors each define a Seifert bundle:
\begin{itemize}
\item \emph{Electroweak} (Paper~IV, Theorem~\ref*{IV-thm:su2}):
  the $\mathrm{SU}(2)_1$ sector operates at
  $N_{\mathrm{EW}} = 4$, defining a line bundle $L_{\mathrm{EW}}$
  with $c_1(L_{\mathrm{EW}}) = 4$.
\item \emph{Gravity/color} (Paper~IV, Theorem~\ref*{IV-thm:su3}):
  the graviton sector operates at
  $N_{\mathrm{grav}} = 7$, defining a line bundle $L_{\mathrm{grav}}$
  with $c_1(L_{\mathrm{grav}}) = 7$.
\end{itemize}
\paragraph{Why $N = 11$ and not another value.}
The cosmological sector must combine the two \emph{gauge}
sectors ($N = 4$ and $N = 7$) on a single fiber.
Other polygon numbers ($N = 3, 5, 6$) carry \emph{matter}
modes but no gauge content (they have $j \ne 1, 2$
at their critical modes).
The selection criteria are:
theorems (KK flux additivity, Thurston uniqueness, McKay);
structural requirements (no duplication, $\ell$ determined);
and physical consistency ($j$ integer, gauge content present).
The selection is \textbf{forced by the central charge budget}:
$c_{\mathrm{gauge}} = 2 + 1 + 1 = 4$ accounts for exactly
$\mathrm{SU}(3)_1 \times \mathrm{SU}(2)_1 \times \mathrm{U}(1)_1$
(Paper~IV, no-hidden-sector theorem).
These three gauge factors require the $N = 7$ gravitational
sector and the $N = 4$ electroweak sector (the unique polygons
with $j = 2$ and $j = 1$ respectively;
Derivation~\ref*{IV-thm:su2}
and Proposition~\ref*{III-prop:graviton}).
On a single $S^1$ fiber (Thurston uniqueness),
flux additivity gives $N = 4 + 7 = 11$.
No alternative exists: adding a third sector would exceed
the central charge budget; removing one would eliminate a
derived gauge factor.

\paragraph{Systematic elimination.}
On a single $S^1$ fiber (Thurston uniqueness,
Paper~IV, Proposition~\ref*{IV-prop:spacetime}),
flux additivity gives
$N/2 = N_{\mathrm{EW}}/2 + N_{\mathrm{grav}}/2 = 4/2 + 7/2 = 11/2$.
All alternatives violate either the Pell equation
(non-integer $j$: mathematical) or the single-sector
requirement (each gauge factor once: structural):
\begin{center}\small
\begin{tabular}{lclc}
\toprule
$N$ & Decomposition & Reason excluded & Type \\
\midrule
$4$  & $4$ alone      & $\Lambda_3 = 0$ (gauge threshold) & math \\
$7$  & $7$ alone      & Gives $M_P/v$, not $\Lambda_4$ & structural \\
$8$  & $4+4$          & No graviton ($j{=}2$ only at $N{=}7$) & math \\
$9$  & $4+5$          & $N=5$: $j \approx 1.79$ (non-integer) & math \\
$10$ & $3+7$          & $N=3$: $j=0$ (no gauge content) & math \\
$12$ & $5+7$          & $N=5$: non-integer $j$ & math \\
$14$ & $7+7$          & Two BO potentials $\to$ extra parameter & structural \\
$28$ & $4 \times 7$   & Doubled gauge sectors & structural \\
\bottomrule
\end{tabular}
\end{center}
\noindent
$N = 11$ is the unique combination including each gauge
sector exactly once.
The Frobenius structure confirms:
$N = 11$'s Frobenius ($m \mapsto 3m \bmod 11$, order~$5$)
gives $\mathrm{SU}(5) \supset \mathrm{SU}(3) \times \mathrm{SU}(2)
\times \mathrm{U}(1)$ (Georgi--Glashow embedding).
\end{proof}

\begin{corollary}[Three generations from the Pell equation]
\label{cor:3-gen}
The number of fermion generations equals the number of
palindromic mode pairs at the graviton polygon:
\begin{equation}\label{eq:n-gen}
  n_{\mathrm{gen}} = \frac{N_{\mathrm{grav}} - 1}{2} = \frac{7 - 1}{2} = 3.
\end{equation}
\end{corollary}

\noindent
Each pair $(m, 7{-}m)$ for $m = 1, 2, 3$ is one generation,
ordered by KK mass $\mu_m = |m - N/2|$ (higher~$m$ gives
larger Yukawa; no permutation freedom).
The Galois orbits $O_+ = \{1, 2, 4\}$ and
$O_- = \{3, 5, 6\}$ each contribute one element per pair,
so $n_{\mathrm{gen}} = n_{\mathrm{color}} = \varphi(7)/2 = 3$.
The three generations carry \emph{distinct} Yukawa textures
from $\mathbb{Z}_7$ charge conservation
(\S\ref*{IV-sec:fermions}).

%% ====================================================================
\section{The electroweak hierarchy}
\label{sec:hierarchy}

\begin{derivation}[Electroweak hierarchy prediction from the hybrid instanton]
\label{thm:hierarchy}
Define, from the inputs $N = 7$, $N = 11$, the Bernoulli
constant $b(N)$ (Paper~I), the Havelock coefficient
$f(3,7) = 6$, the fundamental unit
$\varepsilon_7 = 8 + 3\sqrt{7}$ of $\mathbb{Z}[\sqrt{7}]$,
and the Wheeler--DeWitt mass gap $\Delta\varepsilon$
(\S\ref{sec:mass-gap}), the following three numerical
quantities:
\begin{align}
  \text{\emph{Instanton:}}\quad
  & 2 S_{\mathrm{BO}}(7)
  \;:=\; 2\int_0^{\rho^*}\!\sqrt{2 c_7\,|V_7(\rho)|}\,d\rho
  \;=\; 36.548, \label{eq:hier-instanton} \\
  \text{\emph{Mass gap:}}\quad
  & \Delta\varepsilon\;\ln\varepsilon_7
  \;=\; 2.224, \label{eq:hier-gap} \\
  \text{\emph{Gravity prefactor:}}\quad
  & \tfrac{1}{2}\ln\!\bigl(c_{11}/(24\pi^2)\bigr)
  \;=\; -0.313, \label{eq:hier-gravity}
\end{align}
where $S_{\mathrm{BO}}$ is the WKB tunneling action through
the Born--Oppenheimer potential $V_7(\rho) = \ln(2\sinh\rho)
+ b(7) - 6$ with kinetic coefficient $c_7 = 12\,b(7) = 51.57$
and turning point $\rho^* = 1.734$; $c_{11} = 12\,b(11)$; and
$\Delta\varepsilon = 0.8031$ is the numerically computed
ground-state gap of the Wheeler--DeWitt operator defined by
the inputs $b(7), f(3,7)$ alone (\S\ref{sec:mass-gap}; no
observational data enters its computation).

\textbf{Mathematical claim (proved below).}
The three quantities~\eqref{eq:hier-instanton}--\eqref{eq:hier-gravity}
are each derived from the algebraic inputs (no observational
data enters).  Their sum is
\begin{equation}\label{eq:hierarchy}
  \underbrace{2S_{\mathrm{BO}}(7)}_{36.548}
  + \underbrace{\Delta\varepsilon\;\ln\varepsilon_7}_{2.224}
  + \underbrace{\tfrac{1}{2}\ln\!\bigl(c_{11}/(24\pi^2)\bigr)}_{-0.313}
  \;=\; 38.459,
\end{equation}
computed to $50$ decimal places in the companion code by
exact rational arithmetic for the closed-form pieces combined
with high-precision numerical quadrature for
$S_{\mathrm{BO}}(7)$ and diagonalization for
$\Delta\varepsilon$ (\texttt{src/planetary\_polygons/}
\texttt{proofs/pell\_identity\_proof.py}).

\textbf{Empirical comparison (not proved).}
The observed logarithm of the Planck-to-electroweak scale ratio
is $\ln(M_P/v) = 38.442$, with $M_P = 1.221 \times 10^{19}$~GeV
and $v = 246.22$~GeV.  The prediction~\eqref{eq:hierarchy} differs
from the observation by $\Delta_{\mathrm{obs}} = 0.017$ in the
logarithm, a $0.04\%$ relative deviation in the log or a $1.5\%$
deviation in the ratio $M_P/M_{\mathrm{EW}}$.  This agreement is
a numerical match, not a mathematical equality: the
formula~\eqref{eq:hierarchy} is a prediction derived from the
$N = 7$ and $N = 11$ structural inputs, and the $0.04\%$ residual
is the current tension with observation (open to improvement
via higher-order corrections or to falsification by improved
measurements).

\begin{figure}[t]
\centering
\includegraphics[width=0.7\textwidth]{../figures/fig_hierarchy_decomposition.pdf}
\caption{The electroweak hierarchy decomposed into three
independently computed components.
The instanton action $2S_{\mathrm{BO}}(7) = 36.548$ dominates.
The mass gap $\Delta\varepsilon\,\ln\varepsilon_7 = 2.224$
and gravity prefactor $-0.313$ bring the total to $38.459$,
matching the observed $38.442$ (dashed red line) to $0.04\%$.}
\label{fig:hierarchy}
\end{figure}

\end{derivation}

\begin{proof}
Each component is computed independently.

\paragraph{Why $c_7$ for the instanton but $c_{11}$ for the gravity
prefactor.}
The electroweak hierarchy involves the $N = 7$ sector
(the graviton polygon): $S_{\mathrm{BO}}(7)$ and $\varepsilon_7$
are computed from $b(7)$ and $f(3,7)$.
The gravity prefactor $c_{11}/(24\pi^2)$ uses
$c_{11} = 12\,b(11)$ because the Planck mass
$M_P = \sqrt{c_{11}/G}$ is set by the \emph{full} spacetime
geometry ($N = 11$), not the graviton sector alone.
The two sectors decouple: the $N = 7$ BO potential determines
the tunneling rate (which sets $v/M_P$), while $N = 11$
determines the overall gravitational coupling (which sets
$M_P$ in physical units).
This decoupling is guaranteed by the $\mathbb{Z}_{56}$
charge conservation (Paper~IV, Corollary~\ref*{IV-cor:proton-stability}):
the $N = 7$ and $N = 4$ sectors cannot exchange KK quanta.

\paragraph{Component~1 (Instanton).}
The Born--Oppenheimer potential for the $N = 7$ breathing mode is
$V_7(\rho) = \ln(2\sinh\rho) + b(7) - f(3,7)$,
which is negative for $\rho < \rho^* = 1.734$
and positive for $\rho > \rho^*$.
The WKB tunneling action:
$S = \int_0^{\rho^*} \sqrt{2c_7\,|V_7(\rho)|}\;d\rho = 18.274$
(numerical integration, $c_7 = 12\,b(7) = 51.57$).
The value $b(7)$ enters both $V_7$ (through the potential offset)
and $c_7$ (through the kinetic coefficient); these are the
\emph{same} quantity with \emph{different} roles, not
two independent inputs.
The bounce action is $2S = 36.548$.

\paragraph{Component~2 (Mass gap).}
The asymptotic mass gap of the WDW equation is
$\Delta\varepsilon = 0.8031$
(the numerical WDW eigenvalue gap, \S\ref{sec:mass-gap}).
Numerically $\sqrt{2/\pi} = 0.7979$ matches $\Delta\varepsilon$
to $0.65\%$ (residual $0.0052$); this is an empirical match,
not a derived identity ($V_7$ has no interior critical point,
so the standard harmonic reduction does not apply;
see Supplement~S, Observation~\ref*{SP-prop:harmonic-wdw-gap}).
The hierarchy uses the exact numerical value:
$\Delta\varepsilon\,\ln\varepsilon_7 = 0.8031 \times 2.769 = 2.224$.
The Pell identity~\eqref{eq:pell-identity} below uses
$\sqrt{2/\pi}$ (the analytical part that participates
in the algebraic cancellation).

\paragraph{Component~3 (Gravity prefactor).}
The Planck mass is $M_P = \sqrt{c_{11}/(24\pi^2 R\ell)}$
from the KK reduction $G_4 = G_3/(2\pi R)$
with $G_3 = 3\ell/(2c_{11})$ (Brown--Henneaux).
This gives $(1/2)\ln(c_{11}/(24\pi^2)) = -0.313$
with $c_{11} = 12\,b(11) = 126.56$.
The central charge is $c_{11}$ (not $c_7$) because
the Planck mass is determined by the full spacetime geometry,
which has total Chern class~$N = 11$
(Theorem~\ref{thm:n-selection}).

\paragraph{The Pell-unit numerical identity.}
Combining components~1 and~2 with the regulator
$\ln\varepsilon_7 = \operatorname{arccosh}(8) = 2.7687$
yields the numerical identity
\begin{equation}\label{eq:pell-identity}
  2S_{\mathrm{BO}}(7) + \Delta\varepsilon\,\ln\varepsilon_7
  \;=\; 38.772 \;\approx\; 14\,\ln\varepsilon_7 \;=\; 38.761,
\end{equation}
matching to $0.03\%$
(Supplement~S, Observation~\ref*{SP-obs:pell-numerical}).
Here $\ln\varepsilon_7 = \ln(8{+}3\sqrt{7}\,)$ is the
regulator of the real quadratic field $\mathbb{Q}(\sqrt{7})$
arising from the Pell equation $8^2 - 7 \cdot 3^2 = 1$
(Supplement~S, Lemma~\ref*{SP-lem:pell-regulator}).
What is \emph{proved}, \emph{conjectural}, and
\emph{observed} (cf.\ Supplement~S, Remark~\ref*{SP-rem:pell-status}):
\begin{itemize}
\item \emph{Proved (arithmetic).} Pell uniqueness
  (Proposition~\ref{prop:pell-unique}): $N = 7$ is the unique
  positive integer for which $(N{+}1)^2 - N\lfloor(N{-}1)/2\rfloor^2 = 1$.
  Regulator identification $\operatorname{arccosh}(8)
  = \ln\varepsilon_7$ (Supplement~S,
  Lemma~\ref*{SP-lem:pell-regulator}).
\item \emph{Proved (geometric).} The coefficient
  $14 = 7\cdot 2$ in~\eqref{eq:pell-identity} is the length
  of the primitive closed geodesic on $X(7) = \Gamma(7)\backslash
  \mathbf{H}^2$ represented by the Pell matrix
  $\gamma = \bigl(\begin{smallmatrix} 8 & 21 \\ 3 & 8 \end{smallmatrix}\bigr)$:
  the order of $\bar\gamma$ in $\mathrm{PSL}(2,\mathbb{F}_7)$
  is exactly~$7$ (unipotent conjugacy class), and the
  hyperbolic translation length of $\gamma$ is
  $2\ln\varepsilon_7$, so the primitive geodesic on $X(7)$ has
  length $14\ln\varepsilon_7$
  (Supplement~S, Lemma~\ref*{SP-lem:pell-geodesic}).
  \emph{Therefore the integer $14$ is arithmetic, not fitted.}
\item \emph{Conjectural (trace-formula bridge).}
  The identity~\eqref{eq:pell-identity} is interpreted as a
  Born--Oppenheimer/Selberg trace correspondence: the BO
  tunneling action equals, up to a mass-gap correction, the
  length of the primitive Pell geodesic on $X(7)$
  (Supplement~S, Conjecture~\ref*{SP-conj:bo-selberg}).
  A direct operator equivalence
  $H_{\mathrm{BO}} \sim \Delta_{X(7)}|_{\mathbb{Z}_7}$ is
  \emph{false} (Weyl densities disagree); the conjecture
  asserts only a semiclassical correspondence in the spirit
  of the Langlands ``Beyond Endoscopy'' program
  (\citealt{Langlands2004,FrenkelLanglandsNgo2010}).  A proof
  from first principles is an open problem.
\item \emph{Numerical only.}
  The match $\Delta\varepsilon(c_7) \approx \sqrt{2/\pi}$ to
  $0.65\%$ is empirical (Supplement~S,
  Observation~\ref*{SP-prop:harmonic-wdw-gap}; $V_7$ has no
  interior critical point, so the standard harmonic reduction
  does not apply).  The subleading coefficient
  $4/c_7^2$ in Observation~\ref*{SP-cor:pell-three-term} is a fit
  to the residual, not a derived Dunham coefficient.
\end{itemize}
The hierarchy prediction $M_P/M_{\mathrm{EW}} \approx
\exp(38.459)$ depends on the \emph{numerical values} of
$2S_{\mathrm{BO}}(7) = 36.548$ and
$\Delta\varepsilon\,\ln\varepsilon_7 = 2.224$ obtained by
direct WKB integration and WDW diagonalization; it does
not load-bear on the Pell interpretation of the leading
coefficient.

\begin{proposition}[Pell uniqueness at $N = 7$]
\label{prop:pell-unique}
$N = 7$ is the unique positive integer for which
$(N{+}1, \lfloor(N{-}1)/2\rfloor)$ is a Pell pair,
i.e.\ $(N{+}1)^2 - N \lfloor(N{-}1)/2\rfloor^2 = 1$.
\end{proposition}
\begin{proof}
For odd~$N$ with $m^* = (N{-}1)/2$:
$(N{+}1)^2 - N m^{*2}
= (-N^3{+}6N^2{+}7N{+}4)/4$.
Setting this to~$1$: $N(N{-}7)(N{+}1) = 0$.
The unique positive root is $N = 7$.
For even~$N$ with $m^* = N/2 - 1$:
$(N{+}1)^2 - N m^{*2} = (-N^3 + 8N^2 + 4N + 4)/4$.
Setting this to~$1$: $N(N^2 - 8N - 4) = 0$, whose real roots
$N \in \{0, 4 \pm 2\sqrt{5}\}$ contain no positive integers.
(Full details: Supplement~S,
Proposition~\ref*{SP-prop:pell-unique}.)
\end{proof}

The \emph{per-vortex interpretation} of~\eqref{eq:pell-identity}
(speculative but suggestive): the leading $14 = 2N$ at $N = 7$
suggests that each of the $N = 7$ vortices contributes
$2\ln\varepsilon_7$ to the bounce action, reduced by a
per-vortex subleading correction of order $1/N$.
This interpretation is unproven; the specific numerical
rearrangement into $\sqrt{2/\pi}$ and $4/c_7^2$ pieces is the
empirical form of Observation~\ref*{SP-cor:pell-three-term}
(Supplement~S), where those two numerical occurrences are
logically distinct and tied together only by the
identity~\eqref{eq:pell-identity}.
The instanton provides $94.3\%$ of the Pell exponent;
the mass gap provides the remaining $5.7\%$.
Together: $e^{14\,\ln\varepsilon_7} = \varepsilon_7^{14}
= (8{+}3\sqrt{7}\,)^{14} = 6.82 \times 10^{16}$.

\paragraph{Empirical three-term rearrangement.}
Dividing~\eqref{eq:pell-identity} by $\ln\varepsilon_7$
and substituting $\Delta\varepsilon = \sqrt{2/\pi} + \delta$
with numerical $\delta \approx 0.0052$
(Supplement~S, Observation~\ref*{SP-cor:pell-three-term})
gives the empirical form
\[
  \frac{2S_{\mathrm{BO}}(7)}{\ln\varepsilon_7}
  \;\approx\; 14 \,-\, \sqrt{2/\pi} \,-\, \frac{4}{c_7^2},
\]
where the coincidence $\delta\,c_7^2 \approx 4$ is flagged
as an open numerical observation (not derived from Langer or
Dunham theory).  Restoring the gravity prefactor:
\[
  \ln\frac{M_P}{M_{\mathrm{EW}}}
  \;=\; 2S_{\mathrm{BO}}(7) + \Delta\varepsilon\,\ln\varepsilon_7
  + \tfrac{1}{2}\ln\!\bigl(\tfrac{c_{11}}{24\pi^2}\bigr),
\]
with each term computed numerically
($36.548 + 2.224 - 0.313 = 38.459$).
\end{proof}

\noindent
\emph{Why} $\sim\!10^{17}$ and not larger:
$N_{\mathrm{grav}} = 7$ is the \emph{first} non-trivial Pell solution
(Theorem~\ref{thm:n-selection}).
At the next ($N = 41$):
$\varepsilon_7^{82} \sim 10^{99}$, no atoms could form.

\begin{remark}[Prediction of the Higgs VEV]
The hierarchy formula \emph{predicts}
$v = M_P / \exp(38.459) = 242.0$\,GeV
(observed $246.2$\,GeV, $1.7\%$ off).
The residual reflects higher-order semiclassical corrections
beyond the three numerical terms
$2S_{\mathrm{BO}}(7) + \Delta\varepsilon\,\ln\varepsilon_7
+ \tfrac12\ln(c_{11}/(24\pi^2))$ retained here.
\end{remark}

%% ====================================================================
\section{The cosmological constant from the $N=11$ instanton}
\label{sec:cc-instanton}

With the electroweak VEV $v = 242.0$\,GeV determined by the
$N = 7$ hierarchy (Theorem~\ref{thm:hierarchy}), the $N = 11$
tunneling action fixes the curvature radius~$\ell$.
The following two lemmas isolate the subleading corrections
to the tunneling action; the main theorem references them.

\begin{lemma}[One-loop determinant correction is $-\gamma/2$]
\label{lem:one-loop}
The one-loop functional determinant of the tunneling
amplitude contributes $\delta(\ln A) = -\gamma/2$
to the exponent, where $\gamma = 0.5772\ldots$ is the
Euler--Mascheroni constant.
\end{lemma}

\begin{proof}
The tunneling amplitude is
$A = e^{-S}\bigl(\det'(-\partial^2 + V'')/\det(-\partial^2+\omega^2)
\bigr)^{\!-1/2}$,
with $\det'$ the zero-mode-removed functional determinant.
The ratio is computed via spectral zeta regularization:
$\ln(\det'/\det) = -\zeta'_{\mathrm{rel}}(0)$.
The BF threshold at $c = 1$ has spectral flow $\mathrm{SF} = 1$.
\emph{Proof of $\mathrm{SF} = 1$}:
on $\mathbf{H}^2$, the radial effective potential
for conformal dimension~$c$ is $V_{\mathrm{eff}}(\rho)
= c(c{-}1)/\sinh^2\!\rho$.
As $c$ sweeps from $c_{\mathrm{up}} = 1/2$ to
$c_{\mathrm{dn}} = 3/2$ (the BF crossing of Step~1),
the mass-squared $M^2(c) = c(c{-}1)$ passes through zero
exactly once, at $c = 1$.
The crossing is simple: $dM^2/dc = 2c{-}1 = 1 \ne 0$ at $c = 1$.
By the standard spectral flow theorem (Atiyah--Patodi--Singer),
each simple zero crossing contributes $\mathrm{SF} = \pm 1$;
the sign is $\mathrm{sgn}(dM^2/dc) = +1$.
No other eigenvalue crosses zero in the interval
$c \in [1/2, 3/2]$ (the next zero of $c(c{-}1)$ is at
$c = 0$, outside the interval).
Therefore $\mathrm{SF} = 1$ exactly.
The Hurwitz zeta identity
$\zeta'_H(0, a) = \log\Gamma(a) - \tfrac{1}{2}\log(2\pi)$
gives $\partial_a \zeta'_H(0,a)\big|_{a=1} = \psi(1) = -\gamma$.
The threshold contribution to the relative zeta is therefore
$\zeta'_{\mathrm{threshold}}(0) = -\gamma \times \mathrm{SF} = -\gamma$.
The bulk-spectrum contributions are regular at $s = 0$:
the BO potential
$V_{11}(\rho) = \log(2\sinh\rho) + b - f_{m^*}$
is confining ($V_{11} \to \rho$ as $\rho \to \infty$),
so both $H$ and the reference operator~$H_0$ have purely
discrete spectra $\{\lambda_n\}$ and $\{\lambda_n^0\}$
with Weyl asymptotics $\lambda_n \sim n^2$.
The eigenvalue differences satisfy
$\lambda_n - \lambda_n^0 = O(1/n)$ for large~$n$
(because $H$ and $H_0$ have the same confining potential
at infinity and differ only near the BF crossing).
The relative zeta
$\zeta_{\mathrm{rel}}(s) = \sum(\lambda_n^{-s} - \lambda_n^{0,-s})$
therefore converges for $\mathrm{Re}(s) > 1/2$
and continues meromorphically to $s = 0$
with no pole (the Weyl coefficients match).
The derivative $\zeta'_{\mathrm{rel}}(0)$
receives contributions from two sources:
(i)~the threshold crossing ($\mathrm{SF} = 1$),
which gives $-\gamma$ as computed above;
(ii)~the bulk spectral difference.
The heat-kernel expansion of $H = -\tfrac{1}{2c}\partial_\rho^2
+ V(\rho)$ on the half-line gives
$\mathrm{Tr}(e^{-tH} - e^{-tH_0})
= (a_0' - a_0^{0\prime})\,t^{-1/2}
+ (a_1' - a_1^{0\prime})
+ O(t^{1/2})$
as $t \to 0^+$.
The Seeley--DeWitt coefficients are:
$a_0' = (2c)^{1/2}/(4\pi)^{1/2}$
(volume, identical for $H$ and $H_0$ since both
act on the same half-line),
and $a_1'$ is determined by the \emph{regularized}
potential integral
$\int_0^\infty [V(\rho) - V_{\mathrm{asymp}}(\rho)]\,d\rho$,
where $V_{\mathrm{asymp}} = \rho + \ln 2 + b - f$
is the large-$\rho$ asymptote.
Since $V - V_{\mathrm{asymp}}
= \log(2\sinh\rho) - \rho - \ln 2
= \log(1 - e^{-2\rho})$ (independent of $f$) the
regularized integral is $\int_0^\infty\!\log(1{-}e^{-2\rho})\,
d\rho = -\pi^2/12$
(a standard Bose--Einstein integral)
for \emph{both} $H$ and $H_0$.
Equality $a_0' = a_0^{0\prime}$ holds because both $H$ and
$H_0$ act on the same half-line $[0,\infty)$ with the same
boundary conditions at both endpoints: at $\rho = 0$, the
regular (Friedrichs) self-adjoint extension at the
log-singular limit point
(\citealt{ReedSimon1975II}, ch.~X.1;
the Weyl limit-circle classification applies to
$-\partial^2 + \log(2\sinh\rho) + c$ on $[0,\infty)$);
at $\rho = \infty$, the decaying
boundary condition from the confining asymptote
$V \to \rho$.  Equality $a_1' = a_1^{0\prime}$ follows from the
$f$-independent regularized integral $\int_0^\infty\log(1{-}e^{-2\rho})\,d\rho
= -\pi^2/12$, which determines the $O(1)$ heat-kernel
coefficient for \emph{both} potentials.
Together these ensure that $\zeta_{\mathrm{rel}}(s)$ has
no pole at $s = 0$ or $s = 1/2$.
Explicitly: the heat-kernel remainder
$r(t) = \mathrm{Tr}(e^{-tH} - e^{-tH_0})
- (a_0' - a_0^{0\prime})t^{-1/2}
- (a_1' - a_1^{0\prime})$
satisfies $r(t) = O(t^{1/2})$ as $t \to 0^+$
(from the next Seeley--DeWitt coefficient)
and $r(t) = O(e^{-\delta t})$ as $t \to \infty$
(from the spectral gap $\delta > 0$ of both operators
on the confining half-line).
The Mellin transform
$\zeta_{\mathrm{rel}}(s) = \tfrac{1}{\Gamma(s)}
\int_0^\infty t^{s-1}\mathrm{Tr}(e^{-tH}-e^{-tH_0})\,dt$
therefore converges absolutely for $\mathrm{Re}(s) > 1/2$
and, after subtracting the two (vanishing) Seeley--DeWitt
terms, continues analytically to $\mathrm{Re}(s) > -1/2$
with no poles.
In particular, $\zeta'_{\mathrm{rel}}(0)$ receives no bulk
contribution, only the threshold crossing~(i).
The ratio $\det'/\det = e^\gamma$ follows.
With $(\det)^{-1/2}$:
$\delta(\ln A) = -\gamma/2$,
so $\ln(\ell\,v) = S_{\mathrm{BO}} - \gamma/2$.
\end{proof}

\begin{lemma}[Non-adiabatic corrections vanish for Havelock eigenvectors]
\label{lem:nonadiabatic}
For the Havelock bath, the non-adiabatic coupling between
angular modes vanishes identically:
\begin{equation}\label{eq:no-nonadiabatic}
  \frac{\partial \lambda_m}{\partial\rho}
  = C_1'(\rho) \quad\text{for all } m
  \qquad\Longrightarrow\qquad
  \langle m | \partial H/\partial\rho | m' \rangle
  = C_1'(\rho)\,\delta_{mm'}.
\end{equation}
Consequently, the BO potential is exact at all adiabatic orders;
the only corrections to $S_{\mathrm{BO}}$ are the WKB (Dunham)
anharmonicity series in powers of $1/c_{11}$.
\end{lemma}

\begin{proof}
The eigenvectors of the Havelock Hessian are the
$\mathbb{Z}_N$ Fourier modes
$|m\rangle = N^{-1/2}\sum_j \omega^{mj}$, which are
$\rho$-independent (fixed by the group structure).
Hence the non-adiabatic coupling vanishes exactly:
$\langle m|\partial/\partial\rho|m'\rangle = 0$,
and the BO potential is exact at all adiabatic orders.
\end{proof}

\begin{remark}[Precision budget for Lemma~\ref{lem:nonadiabatic}]
The WKB action
$S_{\mathrm{BO}} = \int_0^{\rho^*}\sqrt{2c\,|V(\rho)|}\,d\rho$
factorizes exactly as $\sqrt{c}$ times a $c$-independent
integral (verified numerically to machine precision for
$N = 7, 11$ across $c \in \{c_N, 2c_N, 4c_N\}$; see
\texttt{tests/test\_wkb\_subleading.py}).  The standard
Dunham cubic-anharmonicity correction
\citep{Dunham1932} at the outer turning point~$\rho^*$,
$\delta S_{\mathrm{Dunham}} = -V'''(\rho^*)/(24 c_N\,|V''(\rho^*)|^{3/2})
\approx e^{\rho^*_{11}}/(12c_{11}) \approx 0.057$,
is \emph{not} the relevant correction: it assumes a smooth
potential near the turning point, but the physical quantity
$\ln(\ell v)$ also depends on the prefactor of the WKB
wavefunction, whose subleading behavior is controlled by the
logarithmic singularity of $V(\rho) = \ln(2\sinh\rho) + b - f$
at $\rho = 0$.  This singularity produces a modified Maslov
index $\alpha \approx 0.40$ (Derivation~\ref{thm:cc},
item~(ii)); the first $O(1/c)$ Dunham correction to
$\ln(\ell v)$ is bounded by $|\delta S| \le 1/c_{11}
\approx 0.008$.
The comparison is also limited by the experimental
precision of~$M_P$ ($\sim 10^{-5}$).
\end{remark}

\begin{derivation}[The cosmological constant from the $N = 11$ instanton]
\label{thm:cc}
The $\mathbf{H}^2$ curvature radius~$\ell$ is determined by the
$N = 11$ Born--Oppenheimer tunneling action:
\begin{equation}\label{eq:ell-cc}
  \ell = \frac{1}{v}\,
  \exp\!\Bigl(S_{\mathrm{BO}}(11)
    - \tfrac{\gamma}{2}
    + O(1/c_{11})\Bigr),
\end{equation}
where $S_{\mathrm{BO}}(11) = 102.724$ is the WKB tunneling action of the
$N = 11$ breathing-mode potential, $\gamma = 0.5772\ldots$ is the
Euler--Mascheroni constant (from Green's function regularization),
$c_{11} = 12\,b(11) = 126.56$, and $v = 242.0$\,GeV (from the
$N = 7$ hierarchy, Theorem~\ref{thm:hierarchy}).
This gives:
\begin{equation}\label{eq:Lambda-pred}
  \ell = 2.450 \times 10^{26}\;\mathrm{m},
  \qquad
  \Lambda_4 = \frac{N^2 - 16}{16\,\ell^2}
  = 4.24 \times 10^{-84}\;\mathrm{GeV}^2,
\end{equation}
predicting
$H_0 = \sqrt{\Lambda_4/(3\,\Omega_\Lambda)} = 67.4\;\mathrm{km\,s^{-1}\,Mpc^{-1}}$,
matching Planck 2018 \citep{Planck2018} ($67.4 \pm 0.5$)
to $< 0.1\%$.
One calibration input ($\Delta m^2_{\mathrm{atm}}$) pins
$M_{\mathrm{poly}}$; $\ell$ is then fixed by the $N = 7$
and $N = 11$ instanton actions alone.
\end{derivation}

\begin{proof}
The derivation chain is:
$M_P$ (the one dimensionful input) $\to$
$v = M_P/\exp(\mathcal{H}_7) = 242.0$\,GeV
(from the $N = 7$ hierarchy exponent
$\mathcal{H}_7 = 2S_{\mathrm{BO}}
+ \Delta\varepsilon\,\ln\varepsilon_7
+ \tfrac{1}{2}\ln(c_{11}/(24\pi^2)) = 38.459$,
Theorem~\ref{thm:hierarchy})
$\to$ $\ell = \exp(S_{\mathrm{BO}}(11))/v$
(from the $N = 11$ instanton, below)
$\to$ $H_0 = \sqrt{\Lambda_4/(3\Omega_\Lambda)}$.
Each step uses only the output of the previous step.

\paragraph{Physical mechanism.}
The $\mathbf{H}^2$ curvature radius~$\ell$ is the scale at which
the $N = 11$ cosmological sector undergoes BO tunneling from the
unstable flat-plane configuration ($\rho = 0$) to the stable
polygon at~$\rho^*_{11}$.
In the WDW equation $[-1/(2c)\,\Psi'' + V(\rho)\,\Psi = 0]$,
the WKB wavefunction between the two turning points
is suppressed by $\exp(-S_{\mathrm{BO}})$.
The ratio of the corresponding length scales
(inner turning point $\sim 1/v$, outer $\sim \ell$) is
therefore $\ell\,v \sim \exp(S_{\mathrm{BO}}(11))$, i.e.\
$\ln(\ell\,v) = S_{\mathrm{BO}}(11)$ at leading order.
This is the same semiclassical relation that gives
the EW hierarchy $M_P/v = \exp(2S_{\mathrm{BO}}(7))$.

The bare exponent is the WKB integral
$S_{\mathrm{BO}}(11) = \int_0^{\rho^*_{11}}
  \sqrt{2\,c_{11}\,|V_{11}(\rho)|}\;d\rho = 102.72$
with $V_{11}(\rho) = \ln(2\sinh\rho) + b(11) - f(5,11)$ and
$\rho^*_{11} = 4.45$.
Two subleading corrections enter:
\begin{enumerate}
\item \emph{The Euler--Mascheroni correction $\gamma/2$}
  (Lemma~\ref{lem:one-loop}):
  the one-loop functional determinant gives
  $\delta(\ln A) = -\gamma/2$,
  so $\ln(\ell\,v) = S_{\mathrm{BO}} - \gamma/2$.
\item \emph{Subleading correction from the log singularity.}
  The potential $V(\rho) = \ln(2\sinh\rho) + b{-}f$ has a
  logarithmic singularity at $\rho = 0$:
  \begin{equation}\label{eq:V-near-zero}
    V(\rho) \;=\; \beta_0 \,+\, \ln\rho
    \,+\, O(\rho^2),
    \qquad \rho \to 0,
    \qquad \beta_0 := \ln 2 + (b{-}f).
  \end{equation}
  Standard WKB breaks down at this singularity because
  the local momentum $p(\rho) = \sqrt{2c\,|V|}$ diverges
  while the exact regular solution $\psi(0) = 1$ is smooth.
  The Liouville substitution $\psi = \rho^{1/2}\chi$,
  $t = -\!\ln\rho$, transforms the WDW equation into
  \begin{equation}\label{eq:chi-equation}
    \chi''(t) \;=\; U(t)\,\chi(t),
    \qquad
    U(t) \;=\; \tfrac{1}{4}
      + 2c\,e^{-2t}\,V\!\bigl(e^{-t}\bigr).
  \end{equation}
  At the boundary $t \to \infty$ ($\rho \to 0$),
  $U \to 1/4$ and $\chi \sim e^{t/2}$; the WKB for $\chi$
  is well-defined at both the boundary and the turning point.
  The connection between these two regimes produces a
  \emph{modified Maslov index}: numerical eigenvalue
  analysis (Bohr--Sommerfeld residual at 54 eigenvalues
  for $N = 7$, $c$ up to $100\,c_N$) gives
  \begin{equation}\label{eq:langer-shift}
    \alpha_{\infty}
      \;\approx\; 0.40,
  \end{equation}
  in place of the standard $\alpha = 1/4$ for a
  regular turning point.  This $O(1)$ phase shift
  is a property of the log singularity and enters the
  Bohr--Sommerfeld quantization, not the WKB amplitude
  directly.

  \emph{Note.}
  The standard \citet{Langer1937} centrifugal prescription
  ($\ell(\ell{+}1) \to (\ell{+}\tfrac{1}{2})^2$) does not
  apply to logarithmic singularities: it is designed for
  the $1/r^2$ barrier in 3D radial equations where the
  first-derivative term has already been removed
  (\citealt{BerryMount1972}~\S IV).
  The first $O(1/c)$ Dunham correction to $\ln(\ell\,v)$
  is bounded by $|\delta S| \le 1/c_{11} \approx 0.008$;
  because the leading terms $S_{\mathrm{BO}}$
  and $-\gamma/2$ already match observation to
  $O(10^{-3})$, this correction is below the current
  precision of the hierarchy prediction.
\end{enumerate}
Together: $\ln(\ell\,v)
= 102.724 - 0.289 + O(1/c_{11}) = 102.435 + O(0.008)$.
The value of $\ell$ is computed from
$S_{\mathrm{BO}}(11)$ (numerical WKB integral), $\gamma/2$
(Hurwitz zeta at BF), and $v$ (from the $N = 7$ hierarchy),
with no cosmological input.
(The discrete choice $N = 11$ is fixed by the selection
principle of Theorem~\ref{thm:n-selection}; $H_0$ is then
a zero-parameter output.)

The BO approximation is exact for the Havelock bath
(Lemma~\ref{lem:nonadiabatic}): the only corrections
beyond the leading WKB integral are the Dunham anharmonicity
series in powers of $1/c_{11}$.
\end{proof}

\paragraph{Verification.} Using $v_{\mathrm{obs}} = 246.22$\,GeV
(the measured value, not the predicted $242.0$\,GeV) and
Planck~2018 \citep{Planck2018} $H_0 = 67.4$\,km/s/Mpc,
$\ln(\ell_{\mathrm{obs}}\,v_{\mathrm{obs}}) = 102.434$.
The match $102.435$ vs.\ $102.434$ gives residual $0.001$,
at the $10^{-5}$ level in the exponent,
well within the $O(1/c_{11})$ Dunham bound.

\paragraph{Sensitivity and theoretical error budget.}
Since $H_0 \propto e^{-S_{\mathrm{BO}}}$, a shift $\delta S$
changes $H_0$ by $\delta S \cdot H_0$.
The corrections form a controlled hierarchy:
\begin{center}
\small
\begin{tabular}{lrl}
\toprule
Contribution & Value & Source \\
\midrule
WKB integral $S_{\mathrm{BO}}(11)$ & $102.7241$ &
  Numerical quadrature ($\pm\!3 \times 10^{-12}$) \\
Euler--Mascheroni $-\gamma/2$ & $-0.2886$ &
  Hurwitz zeta at BF threshold \\
Dunham $O(1/c_{11})$ & $|\delta S| \le 0.008$ &
  Bounded; log-singularity Maslov index $\approx 0.40$ \\
\midrule
$\ln(\ell\,v)$ (predicted) & $102.435 \pm 0.008$ & \\
$\ln(\ell_{\mathrm{obs}}\,v_{\mathrm{obs}})$ (Planck $+$ $v_{\mathrm{obs}}$) & $102.434$ & \\
\midrule
Residual & $0.001$ &
  Within the $O(1/c_{11})$ Dunham bound \\
\bottomrule
\end{tabular}
\end{center}
The numerical integration is not the limiting factor:
SciPy adaptive quadrature converges $S_{\mathrm{BO}}$ to
$10^{-10}$ (independent of the number of subintervals).
The theoretical uncertainty is the unresolved $O(1/c_{11})$
Dunham correction, bounded by $0.008$ in the exponent, which
shifts~$H_0$ by at most $\pm 0.5$\,km/s/Mpc, comparable to
the Planck error bar ($\pm 0.5$\,km/s/Mpc).
Resolving the exact $O(1/c)$ coefficient requires the
connection formula for the non-standard Maslov
index~\eqref{eq:langer-shift} at the log singularity,
which lies outside the scope of the classical
\citeauthor{Langer1937}--\citeauthor{BerryMount1972} theory.
The offset $b(N)$ is proved analytically from the Gauss log-sine
product (Paper~I); no renormalization scheme enters.

\begin{remark}[The Pell connection at $N = 11$]
The $N = 11$ tunneling zero $\cosh\rho^*_{11} \approx 43$ satisfies
the Pell equation
$43^2 - 462 \times 2^2 = 1$,
where $462 = 6 \times 7 \times 11 = f(3,7) \times N_{\mathrm{grav}}
\times N_{\mathrm{cosmo}}$ (the graviton Casimir times both polygon
numbers) and $y = 2 = j$ (the graviton spin).
This parallels the $N = 7$ Pell identity
$8^2 - 7 \times 3^2 = 1$ that controls the electroweak hierarchy.
\end{remark}

%% ====================================================================
\section{The cosmological energy budget}
\label{sec:cosmology}

\subsection{Free energy decomposition}
The $N - 2$ angular Havelock modes with $m \ne m^*$
are \emph{frozen}: their eigenvalues $\lambda_m > 0$ at
the BF threshold prevent them from participating in
breathing-mode tunneling, so they contribute only through
their one-loop determinant.
The partition function $Z = e^{-F}$ of the BO
breathing mode then has three contributions:
\begin{align}
  F_{\mathrm{DE}} &= b(N) + \tfrac{N}{4}
    & &\text{(tree-level: vacuum $+$ radion zero-point)}, \label{eq:F-DE}\\
  F_{\mathrm{DM}} &= \tfrac{1}{2}\!\sum_{m \ne m^*}\!
    \ln\bigl|f(m^*) - f(m)\bigr|
    & &\text{(one-loop: frozen angular modes)}, \label{eq:F-DM}\\
  F_{\mathrm{M}}  &= \Delta E
    & &\text{(non-perturbative: WDW mass gap)}. \label{eq:F-M}
\end{align}
The decomposition is \emph{canonical within the Born--Oppenheimer
saddle-point expansion}: each term arises from a distinct order
(classical action, Gaussian fluctuation, tunneling correction),
and the order assignment follows from the WKB structure
(cf.\ \citealt{BenderOrszag1978}, ch.~10).
The \emph{physical labeling}
(DE $\leftrightarrow$ tree, DM $\leftrightarrow$ one-loop,
baryon $\leftrightarrow$ non-perturbative)
is then determined by the equation of state of each term,
which is derived below from its $\rho$-dependence at the
BO saddle:
\begin{itemize}
\item $F_{\mathrm{DE}}$ is $\rho$-independent at the saddle:
  a constant energy density has $\rho + p = 0$ ($w = -1$),
  which is the equation of state of the cosmological constant
  (dark energy).
\item $F_{\mathrm{DM}} = \tfrac{1}{2}\sum\ln|\lambda_m(\rho)|$
  is the one-loop determinant of the frozen angular modes.
  Unlike $F_{\mathrm{DE}}$ (which is $\rho$-independent),
  $F_{\mathrm{DM}}$ depends on the breathing mode~$\rho$
  through $\lambda_m = C_1(\rho) - f(m)$.
  This $\rho$-dependence is the key:
  a $\rho$-independent contribution would be vacuum energy
  ($w = -1$), but a contribution that depends on the
  dynamical variable~$\rho$ (which evolves with cosmic time)
  has zero pressure.
  Explicitly: the stress-energy tensor of the frozen modes is
  $T^{\mu\nu} = (\rho_E + p)\,u^\mu u^\nu + p\,g^{\mu\nu}$
  with energy density $\rho_E = F_{\mathrm{DM}}(\rho)/(a^3 V_0)$
  (where $a$ is the FRW scale factor and
  $V_0$ the comoving volume) and pressure
  $p = -\partial F_{\mathrm{DM}}/\partial V$.
  Since $F_{\mathrm{DM}} = -\tfrac{1}{2}\sum_{m \ne m^*}
  \ln|\lambda_m(\rho)|$ depends on $\rho$ (the breathing mode,
  evolving with cosmic time via the WDW equation)
  but not on the spatial volume~$V$ of the $\mathbf{H}^2$ slice
  ($\lambda_m$ is determined by the polygon geometry,
  not the slice area):
  $p = -\partial F_{\mathrm{DM}}/\partial V = 0$, hence $w = 0$.
  After the $N$-selection phase transition, each frozen mode
  becomes a massive KK particle with
  $M_k = \sqrt{\lambda_k}\,M_{\mathrm{poly}}$,
  produced at rest in the comoving frame
  (zero peculiar velocity: the BO tunneling creates
  particles at the bottom of the angular potential well,
  not with kinetic energy).
  The stress-energy is
  $T_{\mu\nu} = \rho_E\,u_\mu u_\nu$ with $p = 0$
  (cold dark matter): massive, non-relativistic,
  non-interacting (orthogonal $\mathbb{Z}_N$ Fourier modes
  have zero overlap by character orthogonality).
\item $F_{\mathrm{M}}$ is the non-perturbative
  WDW mass gap (the tunneling amplitude), localized and
  thermalized below $T_{\mathrm{EW}}$: baryonic matter ($w = 0$).
\end{itemize}
Within the Born--Oppenheimer decomposition, the
equation-of-state assignment is determined by the
$\rho$-dependence at the saddle: constant ($w = {-}1$),
frozen massive ($w = 0$), localized tunneling ($w = 0$).
Each assignment follows from the derivation above
(stress-energy computation from the saddle-point action),
not from a free choice of labeling.
The radion contribution $N/4$ in~\eqref{eq:F-DE} is the zero-point
energy of the $S^1$ fiber modulus, stabilized at unit radius by
the flux~$N/2$: the radion mass is $M_{\mathrm{rad}} = N/2$,
giving $(1/2)M_{\mathrm{rad}} = N/4$.
This is determined by the geometry.
The eigenvalues $\lambda_m(\rho^*) = f(m^*) - f(m)$
(where $f(m) = m(N{-}m)/2$) are the Havelock eigenvalues evaluated at
the BF threshold $\rho = \rho^*$, where the critical mode
$\lambda_{m^*} = 0$ and all remaining modes are positive.
The sum excludes both critical modes
($m = m^*$ and $m = N - m^*$ when $N$ is odd).
For $N = 11$: $f(m^*) = 15$, and $F_{\mathrm{DM}} = 5.193$.

\begin{proposition}[Cosmological energy budget at $N = 11$]
\label{prop:cosmology}
At $N = 11$, with the resummed mass gap~\eqref{eq:mass-gap-resum}
and one-loop corrections (\S\ref*{IV-sec:loops}):
\begin{center}
\small
\begin{tabular}{lccc}
\toprule
Component & Prediction & Planck 2018 \citep{Planck2018} & Tension \\
\midrule
Dark energy $\Omega_\Lambda$ & $68.3\%$ & $68.5 \pm 0.7\%$ & $0.3\sigma$ \\
Dark matter $\Omega_{\mathrm{DM}}$ & $26.7\%$ & $26.6 \pm 0.7\%$ & $0.1\sigma$ \\
$\Omega_\Lambda / \Omega_{\mathrm{DM}}$ & $2.56$ & $2.57$ & $0.4\%$ \\
Baryonic matter $\Omega_b$ & $4.97\%$ & $4.93 \pm 0.09\%$ & $0.4\sigma$ \\
\bottomrule
\end{tabular}
\end{center}
\noindent
The three predicted fractions sum to $99.97\%$ ($68.3 + 26.7 + 4.97$);
the remaining $0.03\%$ is the radiation and neutrino fraction
($\Omega_r \sim 10^{-4}$, $\Omega_\nu \sim 10^{-3}$),
not shown in this table, which brings the total to exactly $1$.
\noindent
\begin{figure}[t]
\centering
\includegraphics[width=0.85\textwidth]{../figures/fig_energy_budget.pdf}
\caption{Cosmological energy budget: theory ($N = 11$) versus
Planck 2018. Dark energy, dark matter, and baryonic fractions
all within $1\sigma$.}
\label{fig:energy-budget}
\end{figure}

The dark energy and dark matter fractions are within $1\sigma$ of the
Planck measurements \citep{Planck2018}.
The $\Omega_\Lambda/\Omega_{\mathrm{DM}}$ ratio matches to~$0.4\%$.
The tree-level baryon fraction ($\Delta E = 0.81$) has $7.9\sigma$ tension.
The BF-crossing instanton (the same one giving
$\delta_{\mathrm{CKM}}$ and the mass hierarchy,
\S\ref*{IV-sec:mass-hierarchy}) corrects the mass gap.
The one-instanton contribution:
\begin{equation}\label{eq:mass-gap-inst}
  \Delta E_{\mathrm{1\text{-}inst}}
  = \Delta E_0 \times \bigl(1 + K / n_{\mathrm{gen}}\bigr)
  = 0.81 \times (1 + 0.548/3) = 0.96,
\end{equation}
where $K = e^{-2\pi k_{\mathrm{frac}}} = 0.548$ and
$1/n_{\mathrm{gen}} = 1/3$ counts the BF-crossing generation.
This overshoots ($\Omega_b = 5.1\%$, $1.9\sigma$ above Planck)
because it is the first term of an alternating series.
The $n$-instanton--anti-instanton pairs interfere destructively.
The alternating sign comes from the Reidemeister torsion
of the CS partition function on the Seifert manifold:
for the $n$-th flat $\mathrm{SU}(3)$ connection,
$\tau_n = (-1)^n\,|\tau_n|$
(\citealt{Jeffrey1992};
\citealt{Rozansky1995}).
The sign alternation $(-1)^n$ follows from the
spectral flow of the gauge-covariant Dirac operator:
adjacent flat connections on the Seifert manifold
have spectral flow $\pm 1$
(one eigenvalue crossing zero),
shifting the eta invariant by~$\pm 1$
and flipping the sign of the analytic torsion.
The $n$-instanton amplitude scales as $K^{2n}$ (not
$K^{2n}$ times an $n$-dependent prefactor) because
the CS action on the Seifert manifold is additive
across flat connections: the $n$-th flat connection
has Chern--Simons invariant $\mathrm{CS}_n = n \times
\mathrm{CS}_1$ (mod~$\mathbb{Z}$), so the tunneling
exponent is $e^{-2\pi n \,k_{\mathrm{frac}}} = K^n$.
The instanton--anti-instanton pair has amplitude $K^n
\times K^n = K^{2n}$.
The determinantal prefactors (one-loop around each
saddle) contribute the Reidemeister torsion $|\tau_n|$,
which for the Seifert manifold $\mathbf{H}^2 \times_N
S^1$ is $|\tau_n| = |\tau_1|^n$ by the multiplicative
property of analytic torsion on product connections
\citep{Rozansky1995}.
Combined with the alternating sign: the $n$-th term
is $(-1)^n |\tau_1|^n K^{2n} = (-1)^n (|\tau_1| K^2)^n$.
Since $|\tau_1| \approx 1$ (the torsion of the trivial
connection is normalized) and $K = 0.548$:
the effective expansion parameter is $K^2 = 0.300$.
This gives the geometric resummation
\begin{equation}\label{eq:mass-gap-resum}
  \Delta E_{\mathrm{resum}}
  = \Delta E_0 \times
    \Bigl(1 + \frac{K}{n_{\mathrm{gen}}(1{+}K^2)}\Bigr)
  = 0.81 \times 1.141 = 0.924,
\end{equation}
where the factor $1/(1{+}K^2) = 0.77$ is the standard
resummation of $\sum_{n=0}^\infty (-K^2)^n$.
The series converges by the Leibniz criterion
(alternating, with $|K^2| = 0.300 < 1$).
Truncating at $n = 2$ gives remainder
$|R_2| \le K^6/(1{+}K^2) = 0.021$ ($2.1\%$).
Higher-order coefficients (deviating from the geometric
series) contribute theoretical uncertainty
$O(K^4) \approx 9\%$; the match to Planck
($0.8\sigma$) is within this uncertainty.
No new parameter enters.
The instanton resummation raises $\Delta E$ from
$0.815$ to $0.924$, increasing $\Omega_b$ from the
tree-level $4.22\%$ to $4.67\%$.
The one-loop corrections (\S\ref*{IV-sec:loops})
add a further $+0.30$\,pp, giving
$\Omega_b = 4.97\%$ ($0.4\sigma$ from Planck).
\end{proposition}

The $N = 11$ budget matches Planck \citep{Planck2018} to $< 1\sigma$ for
$\Omega_\Lambda$ and $\Omega_{\mathrm{DM}}$
(Proposition~\ref{prop:cosmology}).
The $\Omega_\Lambda / \Omega_{\mathrm{DM}}$ ratio $= 2.56$
(Planck: $2.57 \pm 0.05$).

\subsection{The dark matter spectrum}

The frozen Havelock modes are the dark matter candidates:
\emph{massive} (from the Casimir gaps), \emph{stable}
($\mathbb{Z}_N$-protected), \emph{non-luminous} (different
KK charge from the Higgs), and \emph{cold} (zero kinetic energy
at threshold).
At $N = 11$, the Casimir gaps are $\lambda_m = 1, 3, 6, 10$
(the \emph{triangular numbers} $T_1, T_2, T_3, T_4$),
giving four dark matter species with mass ratios
\begin{equation}\label{eq:dm-masses}
  M_1 : M_2 : M_3 : M_4
  = 1 : \sqrt{3} : \sqrt{6} : \sqrt{10}.
\end{equation}
There are $8$ dark matter particles in total
($4$ frozen pairs $\times$ degeneracy~$2$).
Their absolute masses are $M_k = \sqrt{\lambda_k} \times M_{\mathrm{poly}}$:
$50$--$160$\,TeV (at $M_{\mathrm{poly}} = 50$\,TeV) or
$300$--$950$\,TeV (at $M_{\mathrm{poly}} = 300$\,TeV).
The lightest species is topologically stable
($\mathbb{Z}_{11}$ charge conservation, cf.\ proton stability
from $\mathbb{Z}_{56}$ in \S\ref*{IV-sec:consistency}).

\paragraph{Production mechanism.}
The frozen modes are \emph{not} thermally produced:
$\sigma \sim 10^{-103}\;\mathrm{cm}^2$ (see
\S\ref*{IV-sec:consistency}) ensures they were never in
thermal equilibrium with the SM bath.
The production mechanism is gravitational particle creation
at the $N$-selection phase transition
($T_N \approx 300$\,TeV, \S\ref{sec:gw-signals}).
During this first-order transition, the breathing mode
$\rho$ tunnels through the BO barrier (bounce action
$2S_{\mathrm{BO}}(11) = 204.5$).
The angular modes, frozen at the saddle with
$\varepsilon_m = 0$ (zero kinetic energy),
become massive KK particles in the broken phase.
The Feynman--Kac correspondence
(Paper~III, Theorem~\ref{III-thm:qm-from-polygons})
maps the quantum ground-state probability $|\Psi_0(\rho)|^2$
to the classical Boltzmann weight $e^{-\beta V}/Z$,
and the one-loop determinant
$Z_{\mathrm{frozen}} = \prod_{m \ne m^*}|\lambda_m|^{-1/2}$
evaluated at the tunneling threshold gives the
number density of each species:
$n_k \propto |\lambda_k|^{-1/2}$.
The equal energy density follows from \emph{equipartition}:
each frozen angular mode~$m$ is a harmonic oscillator
with Hamiltonian $H_m = \tfrac{1}{2}\lambda_m|\varepsilon_m|^2$
at the FK temperature $T = 1$.
By the equipartition theorem:
$\langle H_m\rangle
= \tfrac{1}{2}\lambda_m \times 1/\lambda_m = \tfrac{1}{2}$,
\emph{independent of~$\lambda_m$}.
Therefore $\rho_k = M_k\,n_k
= (\text{modes per species}) \times \tfrac{1}{2} = 1$
for all four species, not a numerical coincidence but a
consequence of classical statistical mechanics.
All four species have \emph{equal energy density},
giving
$\Omega_{\mathrm{DM}} = F_{\mathrm{DM}} / F_{\mathrm{total}}$
(eq.~\ref{eq:F-DM}).
This is the gravitational analogue of Schwinger pair
production: the time-dependent BO potential creates
particles from the vacuum at the phase transition,
with the spectrum set by the eigenvalues, not by
thermal equilibrium.

\paragraph{Clustering.}
After production, the frozen modes are at rest in the
comoving frame (zero peculiar velocity) and cluster
gravitationally like standard cold dark matter:
rotation curves, bullet cluster dynamics, gravitational
lensing, and CMB acoustic peaks are all reproduced by
a CDM component at
$\Omega_{\mathrm{DM}} = 26.7\%$ with masses
$50$--$950$\,TeV.
The $\mathbb{Z}_{11}$ stability ensures that the
lightest frozen mode is absolutely stable (analogous to
$R$-parity in supersymmetric models).

\subsection{The mass gap and baryon fraction}
\label{sec:mass-gap}

The Planck value $\Omega_b/\Omega_m = 15.6\%$ requires
$F_M = 0.96$, above the WDW estimate $\Delta E = 0.815$
(the instanton correction~\eqref{eq:mass-gap-resum} closes
this gap).

The WDW mass gap, computed via non-uniform FEM mesh
(quadratically concentrated near $\rho = 0$), gives converged values:
\begin{center}
\small
\begin{tabular}{lcccccc}
\toprule
$N$ & $6$ & $7$ & $8$ & $9$ & $10$ & $11$ \\
\midrule
$\Delta E$ & $0.841$ & $0.832$ & $0.825$ & $0.821$ & $0.817$ & $0.815$ \\
\bottomrule
\end{tabular}
\end{center}
\noindent
In the semiclassical limit $c_N \to \infty$, the WDW equation
reduces (after rescaling $x = \sqrt{2c_N}\,\rho$) to the
\emph{universal log Schr\"odinger equation}
\begin{equation}\label{eq:log-schrodinger}
  \Bigl[-\tfrac{1}{2\alpha}\,\frac{d^2}{dx^2} + \ln x\Bigr]\Psi
  = \varepsilon\,\Psi
  \qquad (x \in (0,\infty),\; \Psi(0) = 0),
\end{equation}
whose eigenvalue gap
$\Delta\varepsilon = \varepsilon_1 - \varepsilon_0$
is \emph{independent of~$\alpha$}: a universal constant
\begin{equation}\label{eq:mass-gap}
  \Delta\varepsilon = 0.80311 \pm 0.00001
  \qquad\text{(shooting method + Richardson extrapolation;
  see Code Availability)}.
\end{equation}
The WKB quantization condition for the log potential
$V(x) = \ln x$ on $(0,\infty)$ with $\Psi(0) = 0$ is
\[
  \int_0^{x_n}\!\sqrt{2\alpha(\varepsilon_n - \ln x)}\,dx
  = \bigl(n + \tfrac{3}{4}\bigr)\pi,
\]
where $x_n = e^{\varepsilon_n}$ is the turning point
(the $3/4$ accounts for the hard wall at $x = 0$
and the soft turning point at $x = x_n$).
Evaluating the Gaussian integral on the left gives
$x_n\sqrt{\pi\alpha/2}$, hence
$\varepsilon_n = \ln[(n{+}3/4)\sqrt{\pi/\alpha}]$
and $\Delta\varepsilon^{\mathrm{WKB}}
= \ln(7/3) = 0.8473$;
this is an upper bound (WKB overestimates low-lying gaps
for concave potentials; \citealt{Langer1937}).
A lower bound $\Delta\varepsilon > \sqrt{2/\pi} = 0.7979$ follows
from the Gaussian approximation to the ground-state distribution.
No closed-form expression for $\Delta\varepsilon$ is known;
$\Delta\varepsilon$ enters the hierarchy at subleading order
($5.8\%$ of $\ln(M_P/v)$), so the bounded uncertainty
$\sqrt{2/\pi} < \Delta\varepsilon < \ln(7/3)$ propagates as
$< 0.4\%$ of the total.
The full WDW gap at finite~$N$ converges to this constant from above:
$\Delta E(7) = 0.84$, $\Delta E(11) = 0.82$,
$\Delta E(\infty) \to 0.8031$.
At $N = 11$: $\Delta E = 0.815$, giving $\Omega_b = 4.22\%$ at tree level.
Including all one-loop corrections (\S\ref*{IV-sec:loops}) and
the neutrino contribution (\S\ref{sec:neutrinos}):
\begin{center}
\small
\begin{tabular}{lcccc}
\toprule
& Tree & Corrected & Planck 2018 \citep{Planck2018} & Tension \\
\midrule
$\Omega_\Lambda$ & $68.9\%$ & $67.9\%$ & $68.5 \pm 0.7\%$ & $0.8\sigma$ \\
$\Omega_{\mathrm{DM}}$ & $26.9\%$ & $27.6\%$ & $26.6 \pm 0.7\%$ & $1.4\sigma$ \\
$\Omega_b + \Omega_\nu$ & $4.22\%$ & $4.53\%$ & $4.93 \pm 0.06\%$ & $6.6\sigma$ \\
\bottomrule
\end{tabular}
\end{center}
\noindent
Corrections: Coleman--Weinberg ($+0.16$ pp),
neutrino seesaw ($+0.14$ pp), CS level shift ($+0.005$ pp).
With the resummation~\eqref{eq:mass-gap-resum}: $0.4\sigma$
(Proposition~\ref{prop:cosmology}).

The residual $0.40$ percentage-point discrepancy
($4.53\%$ vs.\ $4.93\%$) is sensitive to the Hubble constant:
$\Omega_b = \Omega_b h^2 / h^2$.
At the Planck value $h = 0.674$: $\Omega_b = 4.93\%$.
At the SH0ES local value \citep{Riess2022} $h = 0.730$: $\Omega_b = 4.20\%$.
The predicted $4.53\%$ lies between the Planck and SH0ES
values, suggesting that the residual depends on the resolution
of the $H_0$ tension.

\begin{remark}[Sensitivity to $H_0$]
Since $\Omega_b = \Omega_b h^2 / h^2$, the prediction
is sensitive to the Hubble tension: the $0.40$\,pp residual
($4.53\%$ vs.\ $4.93\%$) corresponds to an $8\%$ relative
error on~$\Omega_b$ for a theory with one discrete criterion
and zero continuous parameters.
\end{remark}

%% ====================================================================
\section{Baryogenesis}
\label{sec:baryogenesis}

The Havelock theory satisfies all three Sakharov conditions.

\paragraph{Baryon number violation.}
The $\mathrm{SU}(2)$ Chern--Simons instantons (the same tunneling
events that generate the hierarchy, Theorem~\ref{thm:hierarchy})
change the CS winding number by~$\pm 1$, violating $B{+}L$
while preserving $B{-}L$.

\paragraph{CP violation.}
The CKM phase $\delta_{\mathrm{CKM}} = \arctan(\sqrt{7}) = 69.3^\circ$
(Paper~IV, Derivation~\ref*{IV-thm:ckm-phase})
arises from the $\mathbb{Z}_7$ quadratic Gauss sum.
The CP-violating Jarlskog invariant
$J = 3.29 \times 10^{-5}$ (Paper~IV, NLO) enters the baryon asymmetry
through the formula below.

\paragraph{First-order electroweak phase transition.}
The Born--Oppenheimer potential $V_7(\rho)$ has a barrier
($S_{\mathrm{BO}} = 18.274$, \S\ref*{IV-sec:loops}).
The EWPT order parameter is determined by the ratio of the
tunneling action to the central charge:
\begin{equation}\label{eq:v-over-Tc}
  \frac{v(T_c)}{T_c}
  = \sqrt{\frac{2\,S_{\mathrm{BO}}}{c}}
  = \sqrt{\frac{36.5}{51.6}} = 0.84.
\end{equation}
This uses the zero-temperature BO potential; finite-temperature
corrections shift $v/T_c$ by $O(T_{\mathrm{EW}}^2/M_{\mathrm{poly}}^2)
\approx 4 \times 10^{-6}$, negligible.
The result is a strong first-order transition (the SM gives
a crossover for $m_H = 125$\,GeV).
In the broken phase, sphaleron processes are suppressed by
$e^{-3v/T_c} = e^{-2.52} = 0.08$: $92\%$ of the baryon asymmetry
survives the transition.

\subsection{Gravitational wave signals}
\label{sec:gw-signals}
The first-order EWPT produces gravitational waves through
bubble collisions, sound waves, and turbulence.
The phase-transition strength
$\alpha = (30/\pi^2 g_*) \times 4D\,(v/T_c)^2 = 0.014$
(with $D = 0.17$ the SM thermal-mass coefficient,
$g_* = 106.75$)
and the inverse duration $\beta/H_* \sim 150$
(from the standard nucleation condition
$S_3(T_*)/T_* \approx 140$)
give a peak frequency and amplitude:
\begin{equation}\label{eq:gw-ewpt}
  f_{\mathrm{peak}} \approx 4 \times 10^{-5}\;\mathrm{Hz},
  \qquad
  \Omega_{\mathrm{GW}}\,h^2
  \approx 10^{-11}\text{--}10^{-9},
\end{equation}
within the LISA sensitivity band ($10^{-4}$--$10^{-1}$\,Hz,
$\Omega h^2 \gtrsim 10^{-12}$).
The SM predicts $\Omega_{\mathrm{GW}} = 0$ (crossover, no signal).
A LISA detection of an EWPT gravitational wave signal
at $f \sim 10^{-5}$\,Hz confirms BSM physics;
the specific $v/T_c = 0.84$ from eq.~\eqref{eq:v-over-Tc}
constrains the BSM sector to the BO barrier shape.

\paragraph{$N$-selection gravitational wave signal.}
The cosmological $N$-selection transition (the tunneling from the
symmetric configuration to the $N = 11$ polygon at $\rho^*_{11} = 4.45$)
is also first-order, with order parameter
$v_N/T_N = \sqrt{2\,S_{\mathrm{BO}}(11)/c_{11}}
= \sqrt{204.5/126.6} = 1.27$
(stronger than the EWPT).
At $T_N \sim M_{\mathrm{poly}} \approx 300$\,TeV,
the GW signal peaks at
$f \approx 0.07$\,Hz with
$\Omega_{\mathrm{GW}}\,h^2 \approx 2 \times 10^{-10}$
(decihertz band; DECIGO, BBO).
The framework predicts two GW signals at distinct frequencies:
the EWPT signal in the LISA band ($10^{-5}$\,Hz)
and the $N$-selection signal in the decihertz band
($10^{-1}$\,Hz).

\subsection{Baryon asymmetry}

The baryon asymmetry is determined by the Bernoulli backbone:
\begin{equation}\label{eq:eta-B-result}
  \eta_B = \frac{J \cdot K^{\binom{N-1}{2}}}{N}
  = \frac{J \cdot K^{15}}{7}
  = 5.7 \times 10^{-10},
\end{equation}
where $J = 3.29 \times 10^{-5}$ is the Jarlskog invariant
computed from the NLO CKM matrix
(Paper~IV, Derivation~\ref*{IV-thm:instanton-vub}:
$|V_{ub}| = |V_{ub}|_{\mathrm{tree}} \times
\mathcal{K}^{N-1}(1{+}\mathcal{K}) = 0.0037$),
$K = e^{-2\pi k_{\mathrm{frac}}} = 0.548$
is the instanton fugacity, and
$\binom{N-1}{2} = \binom{6}{2} = 15$ is the number of
independent off-diagonal entries in the $(N{-}1) \times (N{-}1)$
mass matrix of the $\mathbb{Z}_7$ modes.
The observed value is $\eta_B^{\mathrm{obs}} = 6.12 \times 10^{-10}$
\citep{Planck2018}; the prediction matches within $7\%$
(well within the factor-$3$ theoretical uncertainty from
lattice QCD inputs).

Each independent Yukawa coupling involves one orbifold instanton tunnel;
the baryon asymmetry is suppressed by~$K$ per tunnel.
The Jarlskog~$J$ provides the CP violation,
and $1/N$ is the fractional baryon number from the $\mathbb{Z}_N$ orbifold.
All three inputs ($J$, $K$, $N$) are derived from $N = 7$
with no adjustable parameters.

\subsection{The strong CP problem}
\label{sec:strong-cp}

The determinantal phase of the Yukawa product is
\begin{equation}\label{eq:theta-det}
  \arg\!\det(Y_u \cdot Y_d) = \frac{\pi}{N} = \frac{\pi}{7},
\end{equation}
a topological quantity independent of~$\sigma$.
This is nonzero, so the $\mathbb{Z}_7$ structure does not set
$\bar\theta = 0$ automatically.

The resolution is the Nelson--Barr mechanism:
CP violation in the framework is \emph{spontaneous}
(from the Gauss sum amplitude, not from a fundamental $\theta$ term),
so $\theta_{\mathrm{bare}} = 0$.
The physical $\bar\theta$ receives only loop corrections,
suppressed by the fermion mass hierarchy:
\begin{equation}\label{eq:theta-NB}
  \bar\theta \sim \frac{m_u\,m_d}{m_c\,m_s}
  \cdot \frac{\alpha_s}{4\pi} \cdot \frac{\pi}{N}
  \sim 10^{-12}\;\mathrm{rad},
\end{equation}
well below the experimental limit
$\bar\theta < 10^{-10}$~\citep{Abel2020}.
The prediction $\bar\theta \sim 10^{-12}$ is potentially
accessible to next-generation neutron EDM experiments
(sensitivity goal ${\sim}\,10^{-13}$).

%% ====================================================================
\section{Neutrino masses from the seesaw}
\label{sec:neutrinos}

The right-handed neutrino $\nu_R$ is an $\mathrm{SU}(2)$ singlet
with $N = 4$ KK mass $\mu_4 = 3/2$
(vs.\ $\mu_4 = 1/2$ for the doublet $\nu_L$).
Its Dirac KK mass is \emph{derived} from the Pell identity.
In the Seifert framework, $\nu_R$ is a Kaluza--Klein mode on~$S^1$
with a \emph{Dirac} bulk mass; there is no Majorana mass term
and lepton number is exactly conserved.
The ``seesaw'' is a Dirac (type-III) seesaw: the heavy state
is a Dirac KK mode at scale~$M_{\mathrm{KK}}$, not a Majorana state.
\begin{equation}\label{eq:M-R}
  M_{\mathrm{KK}} = M_{\mathrm{poly}} \times \varepsilon_7^{N_{\mathrm{grav}}}
  \times f(m^*, N_{\mathrm{grav}})
  = M_{\mathrm{poly}} \times \varepsilon_7^7 \times 6
  \sim 10^{13}\text{--}10^{14}\;\mathrm{GeV},
\end{equation}
a structural ansatz determined by three physical inputs:
(i)~$\varepsilon_7^7 = \exp(7\ln\varepsilon_7)
\approx \exp(S_{\mathrm{BO}}(7) + \Delta\varepsilon\,\ln\varepsilon_7/2)$,
from the numerical identity~\eqref{eq:pell-identity}
(the $14 = 2N$ coefficient halved): the BO instanton
that generates the electroweak hierarchy also generates the
KK mass, with the WDW mass gap correction replacing the
bounce factor;
(ii)~$f(3,7) = 6 = j(j{+}1)$ is the graviton Casimir, the
heavy KK mode couples through the $N = 7$
critical mode;
(iii)~$M_{\mathrm{poly}}$ is the polygon scale (from
$\sin^2\theta_W$ running).
One calibration input ($\Delta m^2_{\mathrm{atm}}$) pins
$M_{\mathrm{poly}}$; $M_{\mathrm{KK}}$ is then determined by~$N = 7$.

The Dirac seesaw $m_\nu = m_D^2/M_{\mathrm{KK}}$ with Dirac masses
from properly normalized radial profiles on~$\mathbf{H}^2$
(at warp parameter $\rho^* = 1.734$) gives:
\begin{center}
\small
\begin{tabular}{lccccc}
\toprule
Gen & Pair & $\mu_7$ & $m_D$ & $m_\nu^{\mathrm{diag}}$ \\
\midrule
1 & $(1,6)$ & $2$ & $6.1$\,GeV & $0.0005$\,eV \\
2 & $(2,5)$ & $1$ & $32$\,GeV & $0.013$\,eV \\
3 & $(3,4)$ & $0$ & $63$\,GeV & $0.051$\,eV \\
\bottomrule
\end{tabular}
\end{center}
\noindent
The $m_\nu^{\mathrm{diag}}$ column gives the pre-mixing diagonal
entries ($\sum = 0.065$\,eV).
Diagonalizing the full mass matrix (including the
$\mathbb{Z}/3\mathbb{Z}$ Galois-twist off-diagonal entries
discussed below) yields physical eigenvalues summing to
$\Sigma m_\nu = 0.059$\,eV.
\noindent
Predictions ($M_{\mathrm{KK}}$ from eq.~\eqref{eq:M-R}; see below for
parameter status):
\begin{itemize}
\item \emph{Normal hierarchy}: $m_3 \gg m_2 \gg m_1$.\quad\checkmark
\item $\Delta m^2_{\mathrm{atm}} = 2.45 \times 10^{-3}\;\mathrm{eV}^2$
  (this fixes the overall scale $M_{\mathrm{KK}}$; see the parameter
  discussion below).
\item $\Delta m^2_{\mathrm{sol}} = 7.8 \times 10^{-5}\;\mathrm{eV}^2$
  (observed: $7.5 \times 10^{-5}$; $3\%$ match).
  The \emph{ratio} $\Delta m^2_{\mathrm{atm}}/\Delta m^2_{\mathrm{sol}}
  = 31.5$ (observed: $32.5$, $3\%$) is a genuine prediction:
  it is $M_{\mathrm{KK}}$-independent and determined entirely by the Yukawa texture.
  The improvement from $2.3\times$ (diagonal $M_{\mathrm{KK}}$) to $3\%$
  comes from the $\mathbb{Z}/3\mathbb{Z}$ Galois twist.
  The Frobenius $\sigma\colon m \mapsto 2m \bmod 7$
  maps generation~$i$ to generation~$j$ with unit angular overlap
  (the Fourier modes are related by the Galois action).
  The off-diagonal $M_{\mathrm{KK}}$ entry between generations~$i$
  and~$j$ is the WKB tunneling amplitude through the Casimir barrier:
  $(M_{\mathrm{KK}})_{ij} = M_{\mathrm{KK}} \times
  e^{-\sqrt{|\Delta f|}\,\rho^*/N}$,
  where $\Delta f = |f(m_i) - f(m_j)|$ is the Casimir gap
  between the two modes and $\rho^*/N$ is the WKB path length
  (the geodesic radius divided by the orbifold order).
\item $\Sigma m_\nu = 0.059$\,eV
  $< 0.12$\,eV (Planck bound \citep{Planck2018}).\quad\checkmark
\end{itemize}
The neutrino contribution $\Omega_\nu = 0.14\%$ enters
the baryon budget (Proposition~\ref{prop:cosmology}).

%% ====================================================================
\section{Radion inflation}
\label{sec:inflation}

The radion $\sigma$ (the $S^1$ modulus) has effective potential
\begin{equation}\label{eq:V-radion}
  V(\sigma) = \frac{N^2}{8\sigma^2} - \frac{c_{11}}{12\sigma} + \Lambda\,\sigma,
\end{equation}
where $\sigma$ is the $S^1$ radius (not the inverse),
$N^2/(8\sigma^2)$ is the flux energy,
$c_{11}/(12\sigma)$ is the Casimir energy,
and $\Lambda\,\sigma$ is the cosmological term.
The potential is stabilized at $\sigma_{\min} \approx 1.35$
(for $N = 11$) and grows linearly for large~$\sigma$:
$V \approx \Lambda\,\sigma$.

The radion's initial value is set by the $N$-selection
tunneling (\S\ref{sec:n-selection}): the BO breathing mode
tunnels at $\rho^*_{11} = 4.45$, at which the $S^1$ modulus
takes the value $\sigma_{\mathrm{start}} = \rho^*_{11}/\ln\varepsilon_{11}
\approx 3\,\sigma_{\mathrm{infl}}$.
The radion therefore begins above the inflection point
($\sigma_{\mathrm{start}} > \sigma_{\mathrm{infl}}$)
and rolls toward the minimum, the condition required
for inflection-point inflation.
The full potential~\eqref{eq:V-radion} has an \emph{inflection point} at
\begin{equation}\label{eq:sigma-infl}
  \sigma_{\mathrm{infl}} = \frac{3N^2/8}{c_{11}/12} = \frac{9N^2}{2c_{11}}
  \approx 4.3 \quad (N = 11),
\end{equation}
where $V'' = 0$.
For $\sigma > \sigma_{\mathrm{infl}}$: $V'' < 0$ (concave down,
$\eta < 0$, red tilt).

\paragraph{Inflection-point inflation.}
If inflation occurs near the inflection point (rather than in the
linear regime $\sigma \gg \sigma_{\mathrm{infl}}$), the cubic
approximation gives \citep{AllahverdiDuttaMazumdar2008}
\begin{equation}\label{eq:inflation-predictions}
  n_s = 1 - \frac{8}{3\sqrt{3}\,N_e},
  \qquad r \sim 10^{-6}\;\text{(negligible)}.
\end{equation}
Planck~2018 \citep{Planck2018} ($n_s = 0.965 \pm 0.004$) requires $N_e \approx 45$--$55$,
corresponding to $T_{\mathrm{rh}} \sim 10^6$--$10^9$\,GeV.
The tensor-to-scalar ratio $r \sim 10^{-6}$ is far below the
BICEP/Keck bound \citep{BICEP2021} ($r < 0.036$) and the LiteBIRD sensitivity
($\sigma_r \approx 0.001$): \emph{inflection-point inflation predicts
undetectable gravitational waves from the inflationary epoch}.

In the linear regime ($\sigma \gg \sigma_{\mathrm{infl}}$):
$n_s = 1 - 3/(2N_e)$, $r = 4/N_e$.
At $N_e = 60$: $n_s = 0.975$ ($2.4\sigma$ tension with Planck),
$r = 0.067$ (excluded by BICEP/Keck \citep{BICEP2021}).
The inflection-point mechanism resolves this tension.



\section{Dark matter detection}
\label{sec:dm-detection}

The frozen modes interact with SM matter through the massive
graviphoton ($m_A \sim M_{\mathrm{poly}}$), with coupling
$g \sim m_{\mathrm{KK}} \times M_{\mathrm{poly}} / M_P
\sim 10^{-14}$ (gravitational strength).
The spin-independent cross section on a nucleon:
\[
  \sigma_{\mathrm{SI}} \sim \frac{\mu_N^2}{\pi}
  \left(\frac{g_{\mathrm{DM}}\,g_N}{m_A^2}\right)^2
  \sim 10^{-103}\;\mathrm{cm^2},
\]
vastly below the LZ bound \citep{LZ2023} ($9.2 \times 10^{-48}$\,cm$^2$).
\emph{Prediction}: the frozen Havelock dark matter is invisible
to all current and planned direct detection experiments.
Six properties distinguish the frozen-mode dark matter
from generic CDM and from competing models;
each is falsifiable:
\begin{enumerate}
\item \emph{No self-interaction.}
  The frozen angular modes are orthogonal $\mathbb{Z}_N$
  Fourier modes: their overlap integrals vanish by character
  orthogonality.
  The predicted self-interaction cross section is
  $\sigma_{\mathrm{self}}/m = 0$ at tree level.
  If cluster-merger observations
  (Bullet Cluster, Abell~3827, MACS~J0025)
  or dwarf-galaxy cores establish
  $\sigma_{\mathrm{self}}/m > 0.1\;\mathrm{cm^2/g}$,
  the frozen-mode model is falsified.
\item \emph{No annihilation or decay signal.}
  Each frozen mode is stabilized by
  $\mathbb{Z}_{11}$ charge conservation:
  no decay channel exists within the KK spectrum.
  Indirect-detection experiments
  (Fermi-LAT, CTA, IceCube) should see no
  excess gamma rays, neutrinos, or positrons from
  DM annihilation or decay.
  A confirmed DM annihilation signal at any mass
  falsifies the model.
\item \emph{Cold, not warm.}
  The frozen modes are produced at rest
  (zero peculiar velocity at the phase transition).
  The free-streaming length is
  $\lambda_{\mathrm{fs}} \sim v_0/(H\,a)
  \sim 0$ (non-relativistic at production).
  Lyman-$\alpha$ forest observations constraining
  WDM below $m_{\mathrm{WDM}} > 5$\,keV are automatically
  satisfied; a detection of warm-DM suppression
  in the matter power spectrum at $k > 10\;h/$Mpc
  would falsify the model.
\item \emph{Four species, equal energy density.}
  Equipartition gives $\rho_1 = \rho_2 = \rho_3 = \rho_4$
  (eq.~above), with mass ratios
  $1:\sqrt{3}:\sqrt{6}:\sqrt{10}$.
  Multi-component CDM with equal energy fractions
  produces a step-like transfer function
  $T(k)$ with features at $k_i \propto M_i$
  (the free-streaming scale of each species).
  For masses $50$--$950$\,TeV, the steps lie at
  $k \sim 10^8$--$10^{10}\;h/$Mpc, below the Silk scale
  and unresolvable by current surveys.
  Future 21\,cm intensity mapping (SKA, HERA) at
  $z > 30$ may probe this regime.
\item \emph{No collider signal.}
  The lightest frozen mode at $50$--$300$\,TeV
  exceeds the energy reach of the LHC ($14$\,TeV)
  and planned colliders (FCC-hh at $100$\,TeV
  reaches only the lightest species at the lowest
  polygon scale).
  A DM particle discovered at $m < 50$\,TeV
  falsifies the model.
\item \emph{$N_{\mathrm{eff}}$ unchanged.}
  The frozen modes decouple before BBN
  ($T_N \sim 300$\,TeV $\gg T_{\mathrm{BBN}} \sim 1$\,MeV)
  and contribute $\Delta N_{\mathrm{eff}} = 0$
  to the radiation density.
  CMB-S4 constraining $N_{\mathrm{eff}} = 3.044 \pm 0.03$
  tests this directly.
\end{enumerate}


%% ====================================================================
\section{Spatial geometry and the initial state}
\label{sec:geometry-initial}

Three structural consequences of the Seifert geometry
$\mathbf{H}^2 \times_N S^1$.

\paragraph{Spatial flatness.}
The Kaluza--Klein reduction absorbs the $\mathbf{H}^2$ curvature
($K = -1$) into the three-dimensional cosmological constant
$\Lambda_3 = (N^2 - 16)/16$ (eq.~\eqref{eq:lambda}).
The $3{+}1$-dimensional spatial sections are flat:
$\Omega_k = 0$ at tree level, with perturbative corrections
of order $O(1/c^2) \approx 6 \times 10^{-5}$ at $N = 11$, not
a fine-tuning but a consequence of the KK mechanism, which
absorbs the negative internal curvature and the positive $S^1$
flux energy into~$\Lambda_3$, leaving no residual spatial curvature.
Current constraints ($|\Omega_k| < 0.002$, Planck 2018) are
consistent; next-generation surveys (CMB-S4, Euclid) will reach
$|\Omega_k| < 10^{-4}$, entering the regime where the $O(1/c^2)$
correction is in principle detectable.

\paragraph{No initial singularity.}
The WDW wave function (Paper~III, eq.~\eqref{III-eq:wdw}) is
evanescent in the interior ($V < 0$, $\rho < \rho^*$) and
oscillatory in the exterior ($V > 0$, $\rho > \rho^*$).
As $\rho \to 0$, $V(\rho) \to \log(2\rho) + \mathrm{const} \to -\infty$
and $\Psi \to 0$ exponentially.  The WKB action
\begin{equation}\label{eq:wkb-convergence}
  S_{\mathrm{WKB}} = \int_0^{\rho^*}\!\sqrt{2c\,|V(\rho)|}\;d\rho
  \;<\; \infty
\end{equation}
converges (numerically: $5.8$ oscillations for $N = 7$,
$32.7$ for $N = 11$; see Code Availability).
Classical spacetime begins at the turning point
$\rho = \rho^*$, where the $S^1$ fiber becomes a physical
spatial dimension and Lorentzian time emerges from the
WKB phase.
The Born--Oppenheimer tunneling action $2S = 36.548$
(eq.~\eqref{eq:ell-cc}) determines the amplitude.
The wave function is unique
(Remark~\ref{III-rmk:wdw-uniqueness}): the CMS integrability
of the angular sector selects the ground state of the
effective Hamiltonian, bypassing the Hartle--Hawking vs.\
Vilenkin boundary-condition ambiguity.

\paragraph{Thermodynamic closure.}
The Onsager mechanism requires three prerequisites
(Lemma~\ref{III-lem:onsager-regge}): compact phase
space, fixed total energy, and fixed vortex number~$N$.
All three require a closed system, the negative-temperature
states that drive polygon formation do not exist in open systems
where energy can leak out.
The WDW constraint $H\Psi = 0$ is the quantum version of total
energy conservation.
The universe expands eternally ($w = -1$, $\Lambda > 0$) but
exchanges no energy with an environment.
Local spacetime curvature ($R_{\mu\nu\sigma\tau} \ne 0$) is
present through the full Einstein equations
(Paper~III, Theorem~\ref{III-thm:polygon-einstein}
and Paper~IV, Theorem~\ref{IV-thm:nonlinear-einstein});
the spatial flatness is a statement about the global FLRW
average, not about local geometry.

\paragraph{The Penrose past hypothesis.}
The palindromic threshold $\rho^*$ is the $\beta = 0$
surface of the Onsager microcanonical ensemble:
the critical eigenvalue $\lambda_{m^*}(\rho)$ changes
sign at $\rho = \rho^*$ (verified numerically to
$10^{-4}$; Code Availability).
For $\rho < \rho^*$: the critical mode is inverted
($\lambda_{m^*} < 0$), the effective temperature
is negative, and the regular polygon is the
entropy-maximizing macrostate.
At creation ($\rho = \rho^*$), the gravitational
entropy is
$S_{\mathrm{grav}} = S_{\mathrm{1\text{-}loop}}
\approx 1.1$~nats ($1.6$~bits at $N = 7$).
The maximum (Cardy) entropy is
$S_{\mathrm{BH}} = (\pi c/3)\cosh\rho^* \approx 158$~nats.
The ratio
\begin{equation}\label{eq:penrose-gap}
  \frac{S_{\mathrm{BH}}}{S_{\mathrm{1\text{-}loop}}}
  \approx 144
  \qquad (N = 7;\; \text{growing to } 5.5 \times 10^4
  \text{ at } N = 15)
\end{equation}
is the computable Penrose gap.
Both entropies are evaluated at the same configuration
($\rho = \rho^*$): $S_{\mathrm{1\text{-}loop}}$ is the
microcanonical (one-loop) entropy of the polygon modes;
$S_{\mathrm{BH}}$ is the Cardy (canonical) entropy of the
boundary CFT at the Gibbons--Hawking temperature.
Their ratio measures the non-perturbative state-counting
enhancement (Virasoro descendants beyond the Gaussian sector).
The initial low gravitational entropy is not fine-tuned but
follows from the tunneling occurring at the
negative-to-positive temperature boundary.
The Weyl curvature hypothesis \citep{Penrose1979} is
automatic.
At~$\rho^*$ the polygon is regular: all perturbation
amplitudes $\varepsilon_m = 0$, so the
energy-momentum tensor is isotropic
($T_{\mu\nu} \propto g_{\mu\nu}$),
the Einstein equations give
$R_{\mu\nu} \propto g_{\mu\nu}$
(maximally symmetric), and the Weyl tensor
$W_{\mu\nu\rho\sigma} = 0$.
As~$\rho$ increases past~$\rho^*$, the thermal
perturbation amplitudes
$\langle|\varepsilon_m|^2\rangle = 1/(2\lambda_m)$
(equipartition) break the isotropy: the $m$-th
Fourier mode of~$T_{\mu\nu}$ contributes to the
$2^m$-pole component of the Weyl tensor with weight
$m^2(m^2{-}1)$ (the angular Laplacian eigenvalue
of the traceless-transverse part).
The Weyl invariant
$|W|^2 = \sum_{m \ge 2} m^2(m^2{-}1)\,
\lambda_m\,\langle|\varepsilon_m|^2\rangle
= \tfrac{1}{2}\sum_{m \ge 2} m^2(m^2{-}1)\,\lambda_m$
increases monotonically with~$\rho$
(verified for $N = 7$, $9$, $11$; Code Availability).
Four arrows align in the direction of
increasing~$\rho$: expansion, gravitational entropy,
Weyl curvature, and structure formation, all driven
by the positive-temperature second law acting on
the ordered initial state.


%% ====================================================================
\section{Conclusion}
\label{sec:conclusion}

The $N = 11$ polygon, uniquely selected by KK flux additivity
(Theorem~\ref{thm:n-selection}), determines the cosmological
energy budget through the Born--Oppenheimer tunneling action.
The strongest results are the $\ell$-independent ratios:
$\Omega_\Lambda/\Omega_{\mathrm{DM}} = 2.56$
(Planck: $2.57$, $0.4\%$) and
$\Delta m^2_{\mathrm{sol}}/\Delta m^2_{\mathrm{atm}} = 1/31.5$
(observed: $1/32.5$, $3\%$), both free of calibration inputs.
The absolute scale $H_0 = 67.4$\,km/s/Mpc requires one
dimensionful input ($M_P$) and matches Planck to $< 0.1\sigma$.

The baryon asymmetry
$\eta_B = 5.7 \times 10^{-10}$ matches the observed
$6.12 \times 10^{-10}$ to $7\%$, with all inputs
derived from $N = 7$.
Five falsifiable predictions, normal neutrino hierarchy,
$w = -1$ exactly, no $0\nu\beta\beta$ decay,
$\Omega_k = 0$ exactly (\S\ref{sec:geometry-initial}),
and no initial singularity
(replaced by quantum tunneling at $\rho = \rho^*$), are
tested by JUNO, Euclid, nEXO, CMB-S4, and future
quantum-gravity phenomenology respectively.
The full parameter accounting appears in Paper~VII.

\paragraph{Code availability.}
All numerical verifications are in the companion repository
(\url{https://github.com/gspea/planetary-polygons-unified}).

%% ====================================================================
\bibliographystyle{plainnat}
\bibliography{../shared/refs}

\end{document}
